\documentclass[a4paper]{book}
% preamble_build.tex: copy of Maximus's preamble.tex for the DE v3 build.
% The blocks marked [BUILD COPY CHANGE] are the earlier build fixes.
% The blocks marked [V3 ADDITION] are new in this version (example box,
% spacing). Everything else is unchanged from preamble.tex.

% For margins
\usepackage[a4paper, margin = 20mm]{geometry}

\usepackage[Lenny]{fncychap}
% Rejne, Glenn, Lenny

\usepackage{etoc}

% [V3 ADDITION] cancel: used to show the dx's "cancelling" in Section 11.3
\usepackage{cancel}
\renewcommand{\CancelColor}{\color{red}}

% For enumeration & math symbols
\usepackage{enumerate}
\usepackage{amsmath,amsfonts,amssymb,amsthm,graphicx,lipsum}

% For diagrams, plots and boxes
\usepackage{pgfplots}
\pgfplotsset{compat=1.18}
\usepackage{tikz}
\usepackage[most]{tcolorbox}
\tcbuselibrary{skins}
\usepackage{xcolor}
\pgfplotsset{
    standard/.style={clip = false,
    axis line style = thick,
    trig format=rad,
    axis lines = middle,
    every axis x label/.style={at={(current axis.right of origin)},anchor= west},
    every axis y label/.style={at={(current axis.above origin)},anchor=south}
    }
}

% For GC font
\usepackage{fontspec}
% [BUILD COPY CHANGE 1] The TI-84 key font is loaded only if the file exists
% (checked in ./Fonts/... and ../Fonts/...). Otherwise \ti{} falls back to a
% small boxed sans-serif key. Known letters of the TI font are mapped to their
% key names in the fallback: o = [y=], z = [mode], f = [L3].
\newcommand{\tikey}[1]{\tcbox[on line, enhanced, boxsep=0pt, left=2pt, right=2pt,
  top=1pt, bottom=1pt, boxrule=0.5pt, arc=1.5pt, colback=white,
  colframe=black!70]{\sffamily\footnotesize\strut #1}}
\ExplSyntaxOn
\NewDocumentCommand{\tifallback}{m}{
  \str_case:nnF {#1}
    { {o}{\tikey{y=}} {z}{\tikey{mode}} {f}{\tikey{L3}} }
    { \tikey{#1} }
}
\ExplSyntaxOff
\def\tifontcase{0}
\IfFileExists{Fonts/TI-84 Plus CE Key Fonts/TI84PlusCEKeys.ttf}
  {\def\tifontcase{1}}{}
\IfFileExists{../Fonts/TI-84 Plus CE Key Fonts/TI84PlusCEKeys.ttf}
  {\def\tifontcase{2}}{}
\ifcase\tifontcase\relax
  % 0: font not found, use boxed fallback keys
  \newcommand{\ti}[1]{\tifallback{#1}}
\or
  % 1: original location (as in preamble.tex)
  \newfontfamily\tifont[
    Path = {Fonts/TI-84 Plus CE Key Fonts/},
    UprightFont = {TI84PlusCEKeys.ttf}
  ]{}
  \newcommand{\ti}[1]{{\tifont #1}}
\or
  % 2: one level up (when compiling inside draft-ch1/)
  \newfontfamily\tifont[
    Path = {../Fonts/TI-84 Plus CE Key Fonts/},
    UprightFont = {TI84PlusCEKeys.ttf}
  ]{}
  \newcommand{\ti}[1]{{\tifont #1}}
\fi

% [FONT SETUP copied from h2-math-notes/aoi/notes/preamble_build.tex (DI notes)]
% [AoI items 17/18/37] GC DISPLAY FONT. Every GC key, command, menu item and
% screen line goes through \gcfont, so the font can be swapped HERE in one place.
%   \gc{fnInt(}      command / menu item / stored variable in the display font
%   \gckey{2nd}      a calculator key (display font in a small key outline)
%   \gcscreen{...}   a calculator screen (display font on a grey-green LCD panel)
% The font file lives in ../Fonts (shared by notes, slides and worksheet).
\def\gcfontfile{VT323-Regular.ttf}
\IfFileExists{Fonts/\gcfontfile}{%
  \newfontfamily\gcfont[Path=Fonts/, Scale=1.32]{\gcfontfile}}{%
  \newcommand\gcfont{\ttfamily}}
\newcommand{\gc}[1]{\mbox{\gcfont #1}}
\newcommand{\gckey}[1]{\tcbox[on line, enhanced, boxsep=0pt, left=2pt, right=2pt,
  top=0.6pt, bottom=0.2pt, boxrule=0.5pt, arc=1.5pt, colback=black!4,
  colframe=black!65]{\gcfont\strut #1}}
\definecolor{lcd}{RGB}{214,222,204}
\newcommand{\gcscreen}[1]{\tcbox[on line, enhanced, boxsep=1pt, left=4pt, right=4pt,
  top=2pt, bottom=2pt, boxrule=1.2pt, arc=1pt, colback=lcd, colframe=black!70]{%
  \gcfont\begin{tabular}{@{}l@{}}#1\end{tabular}}}
% integral template as shown on the TI-84 Plus CE (MathPrint): \gcint{0}{2}{e^(-X²)}{X}
\newcommand{\gcint}[4]{{\gcfont\raisebox{-0.35ex}{\Large ∫}\hspace{-0.15em}%
  \raisebox{-0.2ex}{\begin{tabular}{@{}l@{}}\footnotesize #2\\[-1.3ex]\footnotesize #1\end{tabular}}%
  \hspace{0.1em}(#3)d#4}}
% old key macro kept as an alias so nothing breaks
\renewcommand{\ti}[1]{\gckey{#1}}

% SI units
\usepackage{siunitx}

% setspace
%\usepackage{setspace}
%\onehalfspacing

% For hyperlinks within and out of document
\usepackage{hyperref}
\hypersetup{
    colorlinks = true,
    linkcolor = black,
    urlcolor = blue,
}   

% For captions for pics
\usepackage{capt-of}


% tcolorbox custom environments
% Definition ========================================================
\newtcolorbox[auto counter,number within=subsection]{definition}[2][drop shadow]{%
enhanced, before skip=5mm, after skip=6mm, % [V3 ADDITION] spacing
colback=blue!5!white,colframe=blue!50!black,fonttitle=\bfseries,
title= Definition~\thetcbcounter. #2,#1}

% Theorem ==========================================================
\newtcolorbox[auto counter,number within=subsection]{theorem}[2][drop shadow]{%
enhanced, before skip=5mm, after skip=6mm, % [V3 ADDITION] spacing
colback=black!5!white,colframe=black,fonttitle=\bfseries,
title= Theorem~\thetcbcounter. #2,#1}

% Caution box ======================================================
\newtcolorbox{caution}[1][]{enhanced,
before skip=5mm,after skip=6mm, % [V3 ADDITION] was 2mm/3mm
boxrule=1pt,left=11mm,right=2mm,top=3mm,bottom=3mm,
colback=yellow!50,
colframe=yellow!20!black,
sharp corners,rounded corners=southeast,arc is angular,arc=3mm,
underlay={%
\path[fill=tcbcolback!80!black] ([yshift=3mm]interior.south east)--++(-0.4,-0.1)--++(0.1,-0.2);
\path[draw=tcbcolframe,shorten <=-0.05mm,shorten >=-0.05mm] ([yshift=3mm]interior.south east)--++(-0.4,-0.1)--++(0.1,-0.2);
\path[fill=yellow!50!black,draw=none] (interior.south west) rectangle node[white]{\Huge\bfseries !} ([xshift=9mm]interior.north west);
},
drop fuzzy shadow,#1}

% Question box ======================================================
\newtcolorbox{question}[1][]{enhanced,
before skip=5mm,after skip=6mm, % [V3 ADDITION] was 2mm/3mm
boxrule=1pt,left=11mm,right=2mm,top=3mm,bottom=3mm,
colback=blue!20!white,
colframe=blue!20!black,
sharp corners,rounded corners=southeast,arc is angular,arc=3mm,
underlay={%
\path[fill=tcbcolback!80!black] ([yshift=3mm]interior.south east)--++(-0.4,-0.1)--++(0.1,-0.2);
\path[draw=tcbcolframe,shorten <=-0.05mm,shorten >=-0.05mm] ([yshift=3mm]interior.south east)--++(-0.4,-0.1)--++(0.1,-0.2);
\path[fill=blue!50!black,draw=none] (interior.south west) rectangle node[white]{\Huge\bfseries ?} ([xshift=9mm]interior.north west);
},
drop fuzzy shadow,#1}

% Information box ======================================================
\newtcolorbox{info}[1][]{enhanced,
before skip=5mm,after skip=6mm, % [V3 ADDITION] was 2mm/3mm
boxrule=1pt,left=11mm,right=2mm,top=3mm,bottom=3mm,
colback=cyan!20!white,
colframe=cyan!20!black,
sharp corners,rounded corners=southeast,arc is angular,arc=3mm,
underlay={%
\path[fill=tcbcolback!80!black] ([yshift=3mm]interior.south east)--++(-0.4,-0.1)--++(0.1,-0.2);
\path[draw=tcbcolframe,shorten <=-0.05mm,shorten >=-0.05mm] ([yshift=3mm]interior.south east)--++(-0.4,-0.1)--++(0.1,-0.2);
\path[fill=cyan!50!black,draw=none] (interior.south west) rectangle node[white]{\Huge\bfseries i} ([xshift=9mm]interior.north west);
},
drop fuzzy shadow,#1}

% Mistake box ======================================================
\newtcolorbox{mistake}[1][]{enhanced,
before skip=5mm,after skip=6mm, % [V3 ADDITION] was 2mm/3mm
boxrule=1pt,left=11mm,right=2mm,top=3mm,bottom=3mm,
colback=red!20!white,
colframe=red!20!black,
sharp corners,rounded corners=southeast,arc is angular,arc=3mm,
underlay={%
\path[fill=tcbcolback!80!black] ([yshift=3mm]interior.south east)--++(-0.4,-0.1)--++(0.1,-0.2);
\path[draw=tcbcolframe,shorten <=-0.05mm,shorten >=-0.05mm] ([yshift=3mm]interior.south east)--++(-0.4,-0.1)--++(0.1,-0.2);
\path[fill=red!50!black,draw=none] (interior.south west) rectangle node[white]{\Huge\bfseries $\times$} ([xshift=9mm]interior.north west);
},
drop fuzzy shadow,#1}



% [V3 ADDITION] Worked example box ==================================
% Own counter, numbered within the subsection (Example 13.4.5.1; boxes before the first subsection are x.y.0.n; was within the section until 11 Nov 2026) (was: Example 11.2.1, 11.2.2, ...).
% For per-subsection numbering (11.2.1.1) change `section' to `subsection'. Blue `question' boxes are kept for thought-provoking questions
% only; worked examples always use this box. Usage:
%   \begin{example} ...question text... \end{example}
%   \solution{...}
\newtcolorbox[auto counter,number within=subsection]{example}[1][]{%
enhanced, breakable=false, before skip=5mm, after skip=4mm,
colback=teal!5!white, colframe=teal!55!black, boxrule=0.8pt,
left=3mm, right=3mm, top=2mm, bottom=2mm,
sharp corners, rounded corners=northwest, arc=3mm,
fonttitle=\bfseries, coltitle=white,
title=Example~\thetcbcounter, attach boxed title to top left={yshift=-2mm, xshift=3mm},
boxed title style={colback=teal!55!black, sharp corners, rounded corners=northwest, arc=2mm},
drop fuzzy shadow, #1}

% [V3 ADDITION] Spacing between sections, subsections, parts, paragraphs =
% (fncychap's chapter heading is untouched; only \section and below.)
\makeatletter
\renewcommand\section{\@startsection{section}{1}{\z@}%
  {-2ex \@plus -.5ex \@minus -.2ex}{3.2ex \@plus .3ex}%
  {\normalfont\LARGE\bfseries\color{blue!50!black}}}
\renewcommand\subsection{\@startsection{subsection}{2}{\z@}%
  {-1.5ex \@plus -.5ex \@minus -.2ex}{2.6ex \@plus .3ex}%
  {\normalfont\Large\bfseries}}
\renewcommand\subsubsection{\@startsection{subsubsection}{3}{\z@}%
  {-4ex \@plus -1ex \@minus -.5ex}{1.6ex \@plus .2ex}%
  {\normalfont\large\bfseries\color{teal!50!black}}}
\makeatother
\setlength{\parindent}{0pt}
\setlength{\parskip}{0.9ex plus 0.3ex minus 0.2ex}

% Headers and Footers
\usepackage{fancyhdr,lastpage}
% [V3 R2] grey left-margin line on the out-of-syllabus pages (item 13)
\usepackage{eso-pic}
\newcommand{\oosmargin}{\AddToShipoutPictureBG{\put(12mm,0){\color{gray!75}\rule{1.6pt}{\paperheight}}}}
\pagestyle{fancy}
\lhead{\small Vectors I: Notes}
% [BUILD COPY CHANGE 2] Header bug fix: \chaptername is the word "Chapter",
% so the original header printed "Chapter 1: Chapter". \chaptermark now stores
% the bare chapter title in \leftmark, and the header uses \leftmark.
\renewcommand{\chaptermark}[1]{\markboth{#1}{}}
\rhead{\small Compiled by Maximus Pung Jun}
% \cfoot{Page \thepage\  of \pageref{LastPage}} % buggy, does not appear on chapter pages.

% [V3 R2] tick and cross markers for mistake boxes (G3)
\newcommand{\wrong}{\textcolor{red!70!black}{\ensuremath{\boldsymbol{\times}}}}
\newcommand{\cmark}{\textcolor{green!50!black}{\ensuremath{\checkmark}}}

\newcommand{\xmark}{\wrong}   % alias: red cross
% [V3 R2] out-of-syllabus helpers (DE notes item 13)
%   \oosmargin    : grey line in the left margin on every following page
%   \oosmarginoff : stop the margin line
%   \oosnote      : grey banner "Out of syllabus. For your interest only."
%   outofsyllabus : environment = new page + margin line + banner; ends with
%                   \clearpage and switches the margin line off again
\newcommand{\oosmarginoff}{\ClearShipoutPictureBG}
\newcommand{\oosnote}[1][Nothing in this section is examined in H2 Mathematics.]{\par\noindent\colorbox{gray!15}{\parbox{\dimexpr\textwidth-2\fboxsep}{\textbf{Out of syllabus. For your interest only.} #1}}\par}
\newenvironment{outofsyllabus}{\clearpage\oosmargin}{\clearpage\oosmarginoff}

% Notation as in Maximus's notes: italic d, display style everywhere (G1).
% Fractions inside exponents stay small, so never use these macros inside ^{}.
\newcommand{\dydx}{\dfrac{dy}{dx}}
\newcommand{\dvdx}{\dfrac{dv}{dx}}
\newcommand{\dudx}{\dfrac{du}{dx}}
\newcommand{\ddydx}{\dfrac{d^{2}y}{dx^{2}}}
\newcommand{\ddt}[1]{\dfrac{d#1}{dt}}
% Small syllabus flag (no longer used in the DE notes, kept for other chapters)
\newcommand{\syllabusflag}[1]{\par\noindent{\footnotesize\color{red!70!black}\textbf{Syllabus note.} #1}\par}


% Teacher's copy toggle ================================================
% [BUILD COPY CHANGE 3] \ifx compares the whole word (the original
% \if\conform test was always false). DE_notes_student.tex predefines
% \conform as HideSolutions.
% A student copy defines \conform as HideSolutions BEFORE loading this preamble.
\providecommand\conform{ShowSolutions}
\newcommand\ShowSolutionsKey{ShowSolutions}
% [V3 ADDITION] In the student copy the hidden solution leaves stretchable
% blank space (\vfill), so the free space on each page is shared out as
% working space below the examples instead of collecting at the bottom.
\newcommand\solution[1]{\ifx\conform\ShowSolutionsKey
  \par\noindent\textbf{Solution.} #1\par\medskip\else\par\vspace{\workspace}\fi}
% [V3 R2] minimum working space left under each example in the student copy
\newlength\workspace\setlength\workspace{3cm plus 1fill}


\newcommand{\NN}{\mathbb{N}}
\newcommand{\ZZ}{\mathbb{Z}}
\newcommand{\QQ}{\mathbb{Q}}
\newcommand{\RR}{\mathbb{R}}
\newcommand{\CC}{\mathbb{C}}


% =====================================================================
\usepackage{pifont}
% preamble_local.tex: extra packages and styles for the Vectors chapter
% ONLY. Loaded after preamble_build.tex (a verbatim copy of the shared
% book preamble from the PnC / DE build, so that it can be re-synced by
% copying it over again).
% =====================================================================
\usepackage{array, booktabs, makecell}
\usetikzlibrary{arrows.meta, calc, positioning, 3d, decorations.pathreplacing, angles, quotes}
% =====================================================================
% vec_common.tex: shared by the notes (via notes/preamble_local.tex) and
% the slides (via slides/slides_local.tex) for the Vectors chapter.
%
%   CHAPTER NUMBER: change the 13 below and rebuild. The notes set the
%   chapter counter from it, and every slide frame title, Example/Quiz
%   number and box number on the slides uses \chapnum.
% =====================================================================
\def\chapnum{13}

% ---- answers -------------------------------------------------------------
% \ans{key}{value} just prints the value. checks/verify.py reads every
% \ans in the notes and slides and checks it with sympy.
\newcommand{\ans}[2]{#2}

% Column vector, the form used in MF27 and in the examples.
% pmatrix cells are text style, so \dfrac alone still looks small and the
% rows collide. Force display style in every cell and open the rows up.
% Each cell gets a strut as tall as its own content plus \colpad above and
% below, so rows of \dfrac never touch, while integer rows stay compact.
\newlength{\colpad}\setlength{\colpad}{0.55ex}
\newcommand{\colcell}[1]{%
  \setbox0=\hbox{$\displaystyle #1$}%
  \rule[-\dimexpr\dp0+\colpad\relax]{0pt}{\dimexpr\ht0+\dp0+2\colpad\relax}%
  \box0}
\newcommand{\col}[3]{%
  \begingroup
  \renewcommand{\arraystretch}{1.35}%
  \begin{pmatrix}
    \colcell{#1}\\
    \colcell{#2}\\
    \colcell{#3}
  \end{pmatrix}%
  \endgroup
}

% \place{Name}{x}{y}{z} defines a TikZ coordinate. The parentheses are
% inside the macro, so TikZ sees them only after expansion, and
% checks/verify.py can read the drawn coordinates back out.
\newcommand{\place}[4]{\coordinate (#1) at (#2,#3,#4);}

% ---- 3D pictures ----------------------------------------------------------
% Oblique axes: x to the right and slightly down, y out of the page,
% z straight up. Use scale to fit the frame.
\tikzset{
  threeD/.style={
    x={(0.78cm,-0.20cm)},
    y={(0.36cm,0.26cm)},
    z={(0cm,0.78cm)},
    line cap=round, line join=round,
    >=Stealth
  },
  varr/.style={->, thick},
  pt/.style={circle, fill=black, inner sep=1.15pt},
  plane/.style={fill=blue!12, draw=blue!45!black, thick, fill opacity=0.85},
  hid/.style={dashed, black!55}
}
% vecfig.tex: shared 3D diagram style (tikz-3dplot) for Vectors I / II notes and slides.
\usepackage{tikz-3dplot}
\tdplotsetmaincoords{62}{55}
\definecolor{vA}{RGB}{31,90,180}   % first vector / line 1
\definecolor{vB}{RGB}{210,90,20}   % second vector / line 2
\definecolor{vC}{RGB}{0,128,110}   % result / third
\definecolor{vN}{RGB}{140,40,150}  % normals
\definecolor{vP}{RGB}{200,215,240} % plane 1
\definecolor{vQ}{RGB}{245,215,190} % plane 2
\definecolor{vR}{RGB}{200,235,215} % plane 3
\tikzset{
  v3/.style={tdplot_main_coords, line cap=round, line join=round, >=Stealth, scale=#1},
  v3/.default=1,
  ax/.style={->, black!50, thin},
  va/.style={->, very thick, vA},
  vb/.style={->, very thick, vB},
  vc/.style={->, very thick, vC},
  vn/.style={->, very thick, vN},
  ln/.style={very thick},
  dsh/.style={dashed, black!55},
  pp/.style={fill=vP, fill opacity=0.55, draw=vA!70!black, thin},
  pq/.style={fill=vQ, fill opacity=0.55, draw=vB!70!black, thin},
  pr/.style={fill=vR, fill opacity=0.55, draw=vC!70!black, thin},
  dot/.style={circle, fill=black, inner sep=1.2pt},
}
% \vaxes{xlen}{ylen}{zlen}
\newcommand{\vaxes}[3]{%
  \draw[ax] (0,0,0) -- (#1,0,0) node[anchor=west]{$x$};
  \draw[ax] (0,0,0) -- (0,#2,0) node[anchor=south west]{$y$};
  \draw[ax] (0,0,0) -- (0,0,#3) node[anchor=south]{$z$};
  \node[anchor=east] at (0,0,0) {\scriptsize $O$};}
% \rightmark{corner}{along1}{along2} small right-angle square
\newcommand{\rtang}[3]{\draw[thin] ($(#1)!0.25cm!(#2)$) -- ($($(#1)!0.25cm!(#2)$)+($(#1)!0.25cm!(#3)$)-(#1)$) -- ($(#1)!0.25cm!(#3)$);}
\newenvironment{vis}{\par\medskip\begin{center}}{\end{center}\medskip}
% \vang{vertex}{p1}{p2}{label}: small angle arc drawn on the page
\newcommand{\vang}[4]{%
  \coordinate (vg1) at ($(#1)!0.55cm!(#2)$); \coordinate (vg2) at ($(#1)!0.55cm!(#3)$);
  \coordinate (vgm) at ($(vg1)!0.5!(vg2)$);
  \draw[vC, thick] (vg1) .. controls ($(#1)!0.72cm!(vgm)$) .. (vg2);
  \node[vC] at ($(#1)!0.95cm!(vgm)$) {#4};}


% A little air between the parts of an example, and between lines of a list.
\setlength{\itemsep}{0.85ex}
\setlength{\parsep}{0.35ex}

% Student copy: minimum working space under each example. 2.2cm (instead of
% the book's 3cm) is meant to keep the student copy on the same page breaks
% as the teacher copy.
\setlength\workspace{2.2cm plus 1fill}

% =====================================================================
% prelim_fonts.tex: fonts of the real 9758 prelim papers (Word + MathType; pdffonts on
% ACJC 2026 P1/P2, ACJC 2025 P1, JPJC 2026, RVHS 2026, CJC/DHS/TJC/NYJC/HCI/EJC 2025):
%   text Times New Roman; cover/footer Arial; RI page footer Calibri (here Carlito);
%   maths = MathType: Times New Roman letters/digits (italic variables; upright digits,
%   function names, d, e, f), SymbolMT for the integral, sigma, pi and = + < >,
%   scripts 7 pt on 12 pt (about 58 %), tight operator spacing.
% Candidate C (teacher-approved 13 Oct 2026): the real Times New Roman and Arial
% (ttf-mscorefonts-installer) + URW Standard Symbols PS (fonts-urw-base35, a SymbolMT
% clone with the same Symbol encoding) for the integral, sigma, pi, = + < >;
% TeX Gyre Termes Math only for what MathType draws itself (radicals, stretchy brackets,
% fraction bars, arrows). House style: upright d (\mathrm{d}x), upright e and f.
% Loaded once by di_style / de_core (HW sheets) and prelim_style (Tests, solutions).
% =====================================================================
\ifdefined\prelimfontsloaded\endinput\fi
\def\prelimfontsloaded{1}
\edef\pfrestoreat{\catcode`\noexpand\@=\the\catcode`\@\relax}\makeatletter
% MathType script sizes: 7 pt superscripts on 12 pt text (58 %), 5 pt second level
\DeclareMathSizes{12}{12}{7}{5}
\DeclareMathSizes{11}{11}{6.5}{5}
\DeclareMathSizes{10.95}{10.95}{6.5}{5}
\DeclareMathSizes{10}{10}{6}{5}
\DeclareMathSizes{9}{9}{5.5}{5}
\DeclareMathSizes{14.4}{14.4}{8.5}{6}
\setmainfont{Times New Roman}
\setsansfont{Arial}
\newfontfamily\coverfont{Arial}
\newfontfamily\calibrifont{Carlito}
\usepackage{unicode-math}
\setmathfont{TeX Gyre Termes Math}
\setmathfont[range={up/{latin,Latin,num}}]{Times New Roman}
\setmathfont[range={it/{latin,Latin}}]{Times New Roman Italic}
\setmathfont[range={bfup/{latin,Latin,num}}]{Times New Roman Bold}
\setmathfont[range={"002B,"003D,"003C,"003E}]{Standard Symbols PS}
\newfontfamily\pf@sym{Standard Symbols PS}
\newfontfamily\pf@symi[FakeSlant=0.21]{Standard Symbols PS}
% MathType-like spacing (tighter than TeX around = + -)
\AtBeginDocument{\thickmuskip=2.5mu plus 1mu\medmuskip=1.5mu plus 1mu\thinmuskip=3mu
% Symbol integral and sigma (18 pt at 12 pt text, as MathType), slanted Symbol pi
 \renewcommand{\int}{\mathop{\vcenter{\hbox{\pf@sym\fontsize{\fpeval{1.5*\f@size}}{0pt}\selectfont\char"F2}}}\nolimits}%
 \renewcommand{\sum}{\mathop{\vcenter{\hbox{\pf@sym\fontsize{\fpeval{1.5*\f@size}}{0pt}\selectfont\char"E5}}}\limits}%
 \renewcommand{\pi}{\mathord{\text{\pf@symi\char"70}}}}
\pfrestoreat
% Times New Roman + Termes Math (summary-sheet / HW house font)
\usepackage{longtable}
\usepackage{enumitem,multicol,needspace}
\raggedbottom
\newcommand{\vv}[1]{\overrightarrow{#1}}
\newcommand{\cv}[2]{\begin{pmatrix}#1\\#2\end{pmatrix}}
\newcommand{\g}[1]{\textcolor{black!50}{#1}}
\tikzset{ax/.style={->,>=Stealth,black!70},
 vA/.style={->,>=Stealth,very thick,blue!70!black},
 vB/.style={->,>=Stealth,very thick,red!70!black},
 vR/.style={->,>=Stealth,very thick,green!45!black},
 comp/.style={densely dotted,thick,black!60},
 dot/.style={circle,fill=black,inner sep=1.1pt}}
% Grid in the SAME coordinate system as the vectors: step=1 (unitless = 1 x/y unit), integer ends.
\newcommand{\gridax}[4]{\draw[black!15,step=1](#1,#3) grid (#2,#4);
 \draw[ax](#1,0)--(#2+0.5,0) node[right,L]{$x$};
 \draw[ax](0,#3)--(0,#4+0.5) node[above,L]{$y$};}
\tikzset{L/.style={font=\scriptsize}}
% Every diagram cell: same scale, same fixed box (5.6cm x 4.2cm), centred.
% Each diagram gets its own unit u so its grid fills the fixed box (equal x/y units).
% Room is left for point labels (W-1.3cm across, H-0.8cm up); text is NOT scaled.
\newcommand{\D}[5]{\parbox[c][\cellh][c]{6.6cm}{\centering
 \pgfmathsetmacro\u{min((6.4-0.9)/(#2-(#1)+0.5),(\cellh/1cm-0.2-0.45)/(#4-(#3)+0.5))}%
 \begin{tikzpicture}[x=\u cm,y=\u cm]\gridax{#1}{#2}{#3}{#4}#5\end{tikzpicture}}}
\newcommand{\E}[1]{\parbox[c]{\linewidth}{\vspace{3pt}$\displaystyle\begin{aligned}#1\end{aligned}$\vspace{3pt}}}
\newcommand{\N}[1]{\parbox[c][\cellh][c]{\linewidth}{\raggedright #1}}
\newlength{\cellh}\setlength{\cellh}{3.2cm}
\newcommand{\cvv}[3]{\begin{pmatrix}#1\\#2\\#3\end{pmatrix}}
% 3D oblique projection: x toward viewer (down-left), y right, z up. No grids.
\tikzset{td/.style={x={(1cm,0cm)},y={(0.45cm,0.35cm)},z={(0cm,1cm)}}}
% \axt{xmax}{ymin}{ymax}{zmax}: axes only
\newcommand{\axt}[4]{\draw[ax](0,0,0)--(#1+0.6,0,0) node[right,L]{$x$};
 \draw[ax](0,#2,0)--(0,#3+0.6,0) node[above right,L]{$y$};
 \draw[ax](0,0,0)--(0,0,#4+0.6) node[above,L]{$z$};}
% \pj{x}{y}{z}: dotted guides from the point down to the xy-plane and onto the axes
\newcommand{\pj}[3]{\draw[comp](#1,#2,#3)--(#1,#2,0)--(#1,0,0);\draw[comp](#1,#2,0)--(0,#2,0);}
% \DT{unit}{xmax}{ymin}{ymax}{zmax}{body}
\newsavebox{\dtbox}\newlength{\dtmax}
\newcommand{\DT}[6]{\sbox{\dtbox}{\begin{tikzpicture}[td,scale=#1]\axt{#2}{#3}{#4}{#5}#6\end{tikzpicture}}%
 \setlength{\dtmax}{\dimexpr\cellh-0.35cm}%
 \typeout{DTSIZE \the\wd\dtbox\space x \the\dimexpr\ht\dtbox+\dp\dtbox}%
 \ifdim\dimexpr\ht\dtbox+\dp\dtbox>\dtmax \typeout{DTOVER}\fi
 \ifdim\wd\dtbox>6.3cm \typeout{DTWIDE}\fi
 \ifdim\dimexpr\ht\dtbox+\dp\dtbox+0.35cm>\cellh\setlength{\dtmax}{\dimexpr\ht\dtbox+\dp\dtbox+0.35cm}\else\setlength{\dtmax}{\cellh}\fi
 \parbox[c][\dtmax][c]{6.6cm}{\centering\usebox{\dtbox}}}
% generic ratio-theorem sketch: {A label}{C label}{B label}{lambda}{mu}{C position fraction}
\newcommand{\ratiosketch}[6]{\begin{tikzpicture}[scale=0.62]
\coordinate (O) at (0,0);\coordinate (A) at (0.6,2.4);\coordinate (B) at (5,1.6);
\coordinate (C) at ($(A)!#6!(B)$);
\draw[thick,black!70](A)--(B);
\draw[vA](O)--(A);\draw[vB](O)--(B);\draw[vR](O)--(C);
\node[dot,label=below left:$O$] at (O){};
\node[dot,label={[L]above left:#1}] at (A){};
\node[dot,label={[L]right:#3}] at (B){};
\node[dot,label={[L]below=2pt:#2}] at (C){};
\draw[decorate,decoration={brace,amplitude=4pt,raise=3pt}](A)--(C) node[midway,above=8pt,sloped]{$#4$};
\draw[decorate,decoration={brace,amplitude=4pt,raise=3pt}](C)--(B) node[midway,above=8pt,sloped]{$#5$};
\end{tikzpicture}}
\begin{document}
\setcounter{chapter}{\numexpr\chapnum-1\relax}
\chapter{Vectors I}
\localtableofcontents
\clearpage
\section{Vectors in Two Dimensions}
In Secondary 4, we have learned vectors in two dimensions.
\[
\mathbf a=\cv{x}{y}
\]
A vector describes a quantity that has both magnitude and direction. Vectors let us describe position, displacement and direction algebraically, so that geometric problems can be solved by calculation. In H2 Mathematics, we extend vectors from two dimensions to three dimensions, and use them to work with lines and planes, angles, distances, intersections and the foot of the perpendicular. Outside mathematics, vectors are used for forces and velocities in physics, in navigation, in computer graphics and 3D modelling, and throughout engineering. Vectors also lie at the heart of mathematics in linear algebra, which has foundational applications in fields such as physics, engineering, computer science, data science, machine learning and artificial intelligence, economics and cryptography.

\medskip

\begin{center}
\begin{tikzpicture}[x=0.75cm,y=0.75cm]
\draw[black!15,step=1](-3,-3) grid (6,3);
\draw[ax](-3,0)--(6.4,0) node[right,font=\scriptsize]{$x$};
\draw[ax](0,-3)--(0,3.4) node[above,font=\scriptsize]{$y$};
\node[dot,label=below left:$O$] at (0,0){};
\draw[comp](0,0)--(2,0)--(2,1);\draw[comp](3,0)--(3,-2);
\draw[vA](0,0)--(2,1);
\node[dot,blue!70!black,label={[blue!70!black]right:$(2,1)$}] at (2,1){};
\node[blue!70!black,above left=0pt] at (1.9,1.0){$\mathbf a=\cv{2}{1}$};
\draw[vB](0,0)--(3,-2);
\node[dot,red!70!black,label={[red!70!black]right:$(3,-2)$}] at (3,-2){};
\node[red!70!black,below left=0pt] at (1.4,-0.95){$\mathbf b=\cv{3}{-2}$};
\end{tikzpicture}
\end{center}

\medskip
\setlength{\LTleft}{\fill}\setlength{\LTright}{\fill}
\begin{center}
{\renewcommand{\arraystretch}{1}\setlength{\tabcolsep}{5pt}
\begin{longtable}{|>{\raggedright\arraybackslash}m{2.9cm}|>{\raggedright\arraybackslash}m{5.0cm}|>{\centering\arraybackslash}m{6.6cm}|}
\hline
\rule{0pt}{2.6ex}\textbf{Name}&\textbf{Equation}&\textbf{Diagram}\\[0.8ex]\hline\endfirsthead
\hline\rule{0pt}{2.6ex}\textbf{Name}&\textbf{Equation}&\textbf{Diagram}\\[0.8ex]\hline\endhead
\N{Negative of $\mathbf a$\newline $-\mathbf a$}&
\E{-\cv{x}{y}&=\cv{-x}{-y}\\ \g{-\mathbf a}&\g{{}=\cv{-2}{-1}}}&
\D{-4}{3}{-2}{2}{
\draw[vA](0,0)--(2,1);\draw[vB](0,0)--(-2,-1);
\node[dot,label={[L]right:$(2,1)$}] at (2,1){};
\node[dot,label={[L]left:$(-2,-1)$}] at (-2,-1){};
\node[L,blue!70!black,above left] at (1,0.5){$\mathbf a$};
\node[L,red!70!black,below] at (-1.2,-0.7){$-\mathbf a$};}\\\hline
\N{$2\mathbf a$}&
\E{k\cv{x}{y}&=\cv{kx}{ky}\\ \g{2\mathbf a}&\g{{}=2\cv{2}{1}}\\&\g{{}=\cv{4}{2}}}&
\D{-1}{6}{-1}{3}{
\draw[vR](0,0)--(4,2);\draw[vA](0,0)--(2,1);
\node[dot,label={[L]right:$(4,2)$}] at (4,2){};
\node[L,blue!70!black,below right] at (1,0.5){$\mathbf a$};
\node[L,green!45!black,above left] at (3,1.5){$2\mathbf a$};}\\\hline
\N{Magnitude $|\mathbf a|$}&
\E{\left|\cv{x}{y}\right|&=\sqrt{x^2+y^2}\\ \g{|\mathbf a|}&\g{{}=\sqrt{2^2+1^2}}\\&\g{{}=\sqrt5}}&
\D{-1}{4}{-1}{2}{
\draw[comp](2,0)--(2,1);\draw[vA](0,0)--(2,1);
\node[dot,label={[L]above right:$(2,1)$}] at (2,1){};
\node[L,below] at (1,0){$2$};\node[L,right] at (2,0.5){$1$};
\node[L,blue!70!black,above left] at (0.9,0.6){$\sqrt5$};}\\\hline
\N{Triangle law}&
\E{\vv{OA}+\vv{AB}&=\vv{OB}\\ \g{\vv{OB}}&\g{{}=\cv{2}{1}+\cv{3}{-2}}\\&\g{{}=\cv{5}{-1}}}&
\D{-1}{7}{-3}{2}{

\draw[vA](0,0)--(2,1);\draw[vB](2,1)--(5,-1);\draw[vR](0,0)--(5,-1);
\node[L,blue!70!black,above left] at (1.0,0.5){$\vv{OA}$};
\node[L,red!70!black,above right] at (3.5,0){$\vv{AB}$};
\node[L,green!45!black,below] at (3.5,-0.75){$\vv{OB}$};
\node[dot,label={[L]below left:$O$}] at (0,0){};
\node[dot,label={[L]above:$A$}] at (2,1){};
\node[dot,label={[L]right:$B(5,-1)$}] at (5,-1){};}\\\hline
\N{$\vv{AB}$}&
\E{\vv{AB}&=\vv{OB}-\vv{OA}\\ \g{\vv{AB}}&\g{{}=\cv{3}{-2}-\cv{2}{1}}\\&\g{{}=\cv{1}{-3}}}&
\D{-1}{7}{-3}{2}{

\draw[vA](2,1)--(0,0);\draw[vB](0,0)--(3,-2);\draw[vR](2,1)--(3,-2);
\node[dot,label={[L]below left:$O$}] at (0,0){};
\node[dot,label={[L]above:$A(2,1)$}] at (2,1){};
\node[dot,label={[L]below:$B(3,-2)$}] at (3,-2){};
\node[L,green!45!black,left] at (2.3,-0.65){$\vv{AB}$};}\\\hline
\N{Parallelogram law\newline $\mathbf a+\mathbf b$}&
\E{\vv{OA}+\vv{OB}&=\vv{OC}\\ \g{\mathbf a+\mathbf b}&\g{{}=\cv{2}{1}+\cv{3}{-2}}\\&\g{{}=\cv{5}{-1}}}&
\D{-1}{7}{-3}{2}{

\draw[comp](2,1)--(5,-1);\draw[comp](3,-2)--(5,-1);\draw[vA](0,0)--(2,1);\draw[vB](0,0)--(3,-2);\draw[vR](0,0)--(5,-1);
\node[dot,label={[L]below left:$O$}] at (0,0){};
\node[dot,label={[L]above:$A$}] at (2,1){};
\node[dot,label={[L]below:$B$}] at (3,-2){};
\node[dot,label={[L]right:$C(5,-1)$}] at (5,-1){};
\node[L,blue!70!black,above left] at (1,0.5){$\mathbf a$};
\node[L,red!70!black,below left] at (1.5,-1){$\mathbf b$};}\\\hline
\N{Linear combination}&
\E{\lambda\mathbf a+\mu\mathbf b&=\cv{\lambda a_1+\mu b_1}{\lambda a_2+\mu b_2}\\ \g{3\mathbf a-2\mathbf b}&\g{{}=3\cv{2}{1}-2\cv{3}{-2}}\\&\g{{}=\cv{6}{3}-\cv{6}{-4}}\\&\g{{}=\cv{0}{7}}}&
\D{-7}{8}{-1}{8}{

\draw[vA](0,0)--(6,3);\draw[vB](6,3)--(0,7);\draw[vR,line width=1.6pt](0,0)--(0,7);
\node[dot,label={[L]left:$(0,7)$}] at (0,7){};
\node[L,blue!70!black,below right] at (3,1.5){$3\mathbf a$};
\node[L,red!70!black,above right] at (3,5){$-2\mathbf b$};
\node[L,green!45!black,right] at (0.1,3.3){$3\mathbf a-2\mathbf b$};}\\\hline
\N{Equal vectors}&
\E{\vv{OA}&=\vv{CD}\\ \g{\vv{OA}}&\g{{}=\cv{2}{1}}\\ \g{\vv{CD}}&\g{{}=\cv{4-2}{3-2}}\\&\g{{}=\cv{2}{1}}}&
\D{-1}{7}{-1}{4}{
\draw[comp](0,0)--(2,0)--(2,1);\draw[comp](2,2)--(4,2)--(4,3);
\draw[vA](0,0)--(2,1);\draw[vA](2,2)--(4,3);
\node[dot,label={[L]below left:$O$}] at (0,0){};
\node[dot,label={[L]right:$A$}] at (2,1){};
\node[dot,label={[L]above:$C(2,2)$}] at (2,2){};
\node[dot,label={[L]above:$D(4,3)$}] at (4,3){};}\\\hline
\N{Collinear vectors}&
\E{\vv{AB}&=k\,\vv{BC}\\&\Rightarrow A,B,C\text{ collinear}\\ \g{\vv{AB}}&\g{{}=\cv{3-1}{3-2}=\cv{2}{1}}\\ \g{\vv{BC}}&\g{{}=\cv{7-3}{5-3}=\cv{4}{2}}\\&\g{{}=2\,\vv{AB}}}&
\D{-2}{9}{-1}{6}{
\draw[comp](0,1.5)--(8,5.5);
\draw[vA](1,2)--(3,3);\draw[vB](3,3)--(7,5);
\node[dot,label={[L]below right:$A(1,2)$}] at (1,2){};
\node[dot,label={[L]below right:$B(3,3)$}] at (3,3){};
\node[dot,label={[L]below right:$C(7,5)$}] at (7,5){};
\node[dot,label={[L]below left:$O$}] at (0,0){};
\node[L,blue!70!black,above left] at (2,2.5){$\vv{AB}$};
\node[L,red!70!black,above left] at (5,4){$\vv{BC}$};}\\\hline
\end{longtable}}
\end{center}


\clearpage
\section{Vectors in Three Dimensions}
In H2 Mathematics, we extend vectors to three dimensions.
\[
\mathbf a=\cvv{x}{y}{z}
\]

\begin{center}\begin{tikzpicture}[td,scale=0.8]
\axt{4}{-3}{3}{3.6}
\node[dot,label=below right:$O$] at (0,0,0){};
\pj{2}{1}{3}\pj{3}{-2}{2}
\draw[vA](0,0,0)--(2,1,3);
\node[dot,blue!70!black,label={[blue!70!black]right:$(2,1,3)$}] at (2,1,3){};
\node[blue!70!black,left] at (1,0.5,1.7){$\mathbf a=\cvv{2}{1}{3}$};
\draw[vB](0,0,0)--(3,-2,2);
\node[dot,red!70!black,label={[red!70!black]above right:$(3,-2,2)$}] at (3,-2,2){};
\node[red!70!black,right] at (4.6,-2,1.6){$\mathbf b=\cvv{3}{-2}{2}$};
\end{tikzpicture}\end{center}

\begin{center}
{\renewcommand{\arraystretch}{1}\setlength{\tabcolsep}{5pt}
\begin{longtable}{|>{\raggedright\arraybackslash}m{2.9cm}|>{\raggedright\arraybackslash}m{5.0cm}|>{\centering\arraybackslash}m{6.6cm}|}
\hline
\rule{0pt}{2.6ex}\textbf{Name}&\textbf{Equation}&\textbf{Diagram}\\[0.8ex]\hline\endfirsthead
\hline\rule{0pt}{2.6ex}\textbf{Name}&\textbf{Equation}&\textbf{Diagram}\\[0.8ex]\hline\endhead
\N{Negative of $\mathbf a$\newline $-\mathbf a$}&
\E{-\cvv{x}{y}{z}&=\cvv{-x}{-y}{-z}\\ \g{-\mathbf a}&\g{{}=\cvv{-2}{-1}{-3}}}&
\DT{0.36}{3}{-3}{3}{3}{
\draw[vA](0,0,0)--(2,1,3);\draw[vB](0,0,0)--(-2,-1,-3);
\node[dot,label={[L]right:$(2,1,3)$}] at (2,1,3){};
\node[dot,label={[L]left:$(-2,-1,-3)$}] at (-2,-1,-3){};
\node[L,blue!70!black,left] at (1,0.5,1.5){$\mathbf a$};
\node[L,red!70!black,right] at (-1,-0.5,-1.5){$-\mathbf a$};}\\\hline
\N{$2\mathbf a$}&
\E{k\cvv{x}{y}{z}&=\cvv{kx}{ky}{kz}\\ \g{2\mathbf a}&\g{{}=2\cvv{2}{1}{3}=\cvv{4}{2}{6}}}&
\DT{0.34}{4}{-3}{6}{6}{
\draw[vR](0,0,0)--(4,2,6);\draw[vA](0,0,0)--(2,1,3);
\node[dot,label={[L]right:$(4,2,6)$}] at (4,2,6){};
\node[L,blue!70!black,left] at (1,0.5,1.5){$\mathbf a$};
\node[L,green!45!black,left] at (3,1.5,4.5){$2\mathbf a$};}\\\hline
\N{Magnitude $|\mathbf a|$}&
\E{\left|\cvv{x}{y}{z}\right|&=\sqrt{x^2+y^2+z^2}\\ \g{|\mathbf a|}&\g{{}=\sqrt{2^2+1^2+3^2}}\\&\g{{}=\sqrt{14}}}&
\DT{0.48}{3}{-3}{3}{3}{\pj{2}{1}{3}
\draw[vA](0,0,0)--(2,1,3);
\node[dot,label={[L]right:$(2,1,3)$}] at (2,1,3){};
\node[L,below] at (1,0,0){$2$};\node[L,below right] at (2,0.5,0){$1$};\node[L,right] at (2,1,1.3){$3$};
\node[L,blue!70!black,left] at (1,0.5,1.6){$\sqrt{14}$};}\\\hline
\N{Triangle law}&
\E{\vv{OA}+\vv{AB}&=\vv{OB}\\ \g{\vv{OB}}&\g{{}=\cvv{2}{1}{3}+\cvv{3}{-2}{2}}\\&\g{{}=\cvv{5}{-1}{5}}}&
\DT{0.36}{5}{-3}{3}{5}{
\draw[vA](0,0,0)--(2,1,3);\draw[vB](2,1,3)--(5,-1,5);\draw[vR](0,0,0)--(5,-1,5);
\node[L,blue!70!black,left] at (1,0.5,1.5){$\vv{OA}$};
\node[L,red!70!black,above left] at (3.5,0,4){$\vv{AB}$};
\node[L,green!45!black,below right] at (2.5,-0.5,2.5){$\vv{OB}$};
\node[dot,label={[L]below right:$O$}] at (0,0,0){};
\node[dot,label={[L]above left:$A$}] at (2,1,3){};
\node[dot,label={[L]right:$B(5,-1,5)$}] at (5,-1,5){};}\\\hline
\N{$\vv{AB}$}&
\E{\vv{AB}&=\vv{OB}-\vv{OA}\\ \g{\vv{AB}}&\g{{}=\cvv{3}{-2}{2}-\cvv{2}{1}{3}}\\&\g{{}=\cvv{1}{-3}{-1}}}&
\DT{0.42}{3}{-3}{3}{3}{
\draw[vA](2,1,3)--(0,0,0);\draw[vB](0,0,0)--(3,-2,2);\draw[vR](2,1,3)--(3,-2,2);
\node[dot,label={[L]below right:$O$}] at (0,0,0){};
\node[dot,label={[L]above:$A(2,1,3)$}] at (2,1,3){};
\node[dot,label={[L]right:$B(3,-2,2)$}] at (3,-2,2){};
\node[L,green!45!black,right] at (2.5,-0.5,2.5){$\vv{AB}$};}\\\hline
\N{Parallelogram law\newline $\mathbf a+\mathbf b$}&
\E{\vv{OA}+\vv{OB}&=\vv{OC}\\ \g{\mathbf a+\mathbf b}&\g{{}=\cvv{2}{1}{3}+\cvv{3}{-2}{2}}\\&\g{{}=\cvv{5}{-1}{5}}}&
\DT{0.5}{5}{-3}{3}{5}{
\draw[comp](2,1,3)--(5,-1,5);\draw[comp](3,-2,2)--(5,-1,5);\draw[vA](0,0,0)--(2,1,3);\draw[vB](0,0,0)--(3,-2,2);\draw[vR](0,0,0)--(5,-1,5);
\node[dot,label={[L]below right:$O$}] at (0,0,0){};
\node[dot,label={[L]above left:$A$}] at (2,1,3){};
\node[dot,label={[L]right:$B$}] at (3,-2,2){};
\node[dot,label={[L]right:$C(5,-1,5)$}] at (5,-1,5){};
\node[L,blue!70!black,left] at (1,0.5,1.5){$\mathbf a$};
\node[L,red!70!black,below right] at (1.5,-1,1){$\mathbf b$};}\\\hline
\N{Linear combination}&
\E{\lambda\mathbf a+\mu\mathbf b&=\cvv{\lambda a_1+\mu b_1}{\lambda a_2+\mu b_2}{\lambda a_3+\mu b_3}\\ \g{3\mathbf a-2\mathbf b}&\g{{}=\cvv{6}{3}{9}-\cvv{6}{-4}{4}}\\&\g{{}=\cvv{0}{7}{5}}}&
\DT{0.27}{6}{-3}{7}{9}{
\draw[vA](0,0,0)--(6,3,9);\draw[vB](6,3,9)--(0,7,5);\draw[vR,line width=1.6pt](0,0,0)--(0,7,5);
\node[dot,label={[L]left:$(0,7,5)$}] at (0,7,5){};
\node[L,blue!70!black,right] at (3,1.5,4.5){$3\mathbf a$};
\node[L,red!70!black,above] at (3,5,7.2){$-2\mathbf b$};
\node[L,green!45!black,left] at (0,3.5,2.5){$3\mathbf a-2\mathbf b$};}\\\hline
\N{Equal vectors}&
\E{\vv{OA}&=\vv{CD}\\ \g{\vv{OA}}&\g{{}=\cvv{2}{1}{3}}\\ \g{\vv{CD}}&\g{{}=\cvv{3-1}{4-3}{4-1}=\cvv{2}{1}{3}}}&
\DT{0.42}{3}{-3}{4}{4}{
\draw[vA](0,0,0)--(2,1,3);\draw[vA](1,3,1)--(3,4,4);
\node[dot,label={[L]below right:$O$}] at (0,0,0){};
\node[dot,label={[L]left:$A$}] at (2,1,3){};
\node[dot,label={[L]right:$C(1,3,1)$}] at (1,3,1){};
\node[dot,label={[L]right:$D(3,4,4)$}] at (3,4,4){};}\\\hline
\N{Collinear vectors}&
\E{\vv{AB}&=k\,\vv{BC}\\&\Rightarrow A,B,C\text{ collinear}\\ \g{\vv{AB}}&\g{{}=\cvv{2-1}{3-2}{3-1}=\cvv{1}{1}{2}}\\ \g{\vv{BC}}&\g{{}=\cvv{4-2}{5-3}{7-3}=\cvv{2}{2}{4}}\\&\g{{}=2\,\vv{AB}}}&
\DT{0.3}{4}{-3}{2}{7}{
\draw[vA](1,2,1)--(2,3,3);\draw[vB](2,3,3)--(4,5,7);
\node[dot,label={[L]right:$A(1,2,1)$}] at (1,2,1){};
\node[dot,label={[L]right:$B(2,3,3)$}] at (2,3,3){};
\node[dot,label={[L]right:$C(4,5,7)$}] at (4,5,7){};
\node[dot,label={[L]below right:$O$}] at (0,0,0){};
\node[L,blue!70!black,left] at (1.5,2.5,2){$\vv{AB}$};
\node[L,red!70!black,left] at (3,4,5){$\vv{BC}$};}\\\hline
\end{longtable}}
\end{center}


\clearpage
\section{Basic Terminology and Notation}
\begin{definition}{{Unit vectors $\mathbf i$, $\mathbf j$, $\mathbf k$}}
\noindent\begin{minipage}[c]{0.58\linewidth}
$\mathbf i$, $\mathbf j$ and $\mathbf k$ are the unit vectors in the positive directions of the $x$-, $y$- and $z$-axes:
\[
\mathbf i=\cvv{1}{0}{0},\qquad \mathbf j=\cvv{0}{1}{0},\qquad \mathbf k=\cvv{0}{0}{1}
\]
Any vector can be written in terms of $\mathbf i$, $\mathbf j$, $\mathbf k$:
\[
\cvv{x}{y}{z}=x\mathbf i+y\mathbf j+z\mathbf k
\]
\end{minipage}\hfill
\begin{minipage}[c]{0.38\linewidth}\centering
\begin{tikzpicture}[td,scale=1.25]
\axt{1.6}{-0.3}{1.8}{1.6}
\node[dot,label={[L]below right:$O$}] at (0,0,0){};
\draw[vA](0,0,0)--(1,0,0) node[midway,below,L]{$\mathbf i$};
\draw[vB](0,0,0)--(0,1,0) node[midway,below right,L]{$\mathbf j$};
\draw[vR](0,0,0)--(0,0,1) node[midway,left,L]{$\mathbf k$};
\end{tikzpicture}
\end{minipage}
\end{definition}

\begin{example}
Write in column vector form: (a) $2\mathbf i+3\mathbf j-4\mathbf k$ \quad (b) $3\mathbf i+2\mathbf k$. \quad (c) Write $\cvv{5}{-1}{0}$ in terms of $\mathbf i,\mathbf j,\mathbf k$.
\end{example}
\solution{
\[
\begin{aligned}
\text{(a)}\quad 2\mathbf i+3\mathbf j-4\mathbf k&=\cvv{2}{3}{-4}
\qquad\qquad
\text{(b)}\quad 3\mathbf i+2\mathbf k=3\mathbf i+0\mathbf j+2\mathbf k=\cvv{3}{0}{2}
\qquad\qquad
\text{(c)}\quad \cvv{5}{-1}{0}=5\mathbf i-\mathbf j
\end{aligned}
\]}

\begin{theorem}{Magnitude of a vector}
\[
\left|\cvv{x}{y}{z}\right|=\sqrt{x^2+y^2+z^2}
\]
\end{theorem}

\begin{example}
Find $\left|\cvv{2}{3}{-4}\right|$.
\end{example}
\solution{\par\noindent\begin{minipage}[c]{0.55\linewidth}
\[
\begin{aligned}
\left|\cvv{2}{3}{-4}\right|&=\sqrt{2^2+3^2+(-4)^2}\\
&=\sqrt{4+9+16}\\
&=\sqrt{29}
\end{aligned}
\]
\end{minipage}\hfill
\begin{minipage}[c]{0.4\linewidth}\centering
\begin{tikzpicture}[td,scale=1.05]
\axt{2.8}{-1}{4}{0.6}
\draw[ax](0,0,0)--(0,0,-4.3);
\node[dot,label={[L]above left:$O$}] at (0,0,0){};
\draw[comp](0,0,0)--(2,0,0)--(2,3,0)--(2,3,-4);
\draw[comp](0,3,0)--(2,3,0);
\node[L,above] at (1,3,0){$2$};
\node[L,below right] at (2,1.5,0){$3$};
\node[L,right] at (2,3,-2){$4$};
\draw[vA](0,0,0)--(2,3,-4);
\node[blue!70!black,above right] at (1,1.5,-2){$\sqrt{29}$};
\node[dot,label={[L]below:$P(2,3,-4)$}] at (2,3,-4){};
\end{tikzpicture}
\end{minipage}}

\begin{theorem}{Parallel vectors}
\[
\mathbf a\parallel\mathbf b\iff\mathbf a=\lambda\mathbf b\ \text{ for some scalar }\lambda\neq0
\quad(\text{alternatively, }\mathbf a\times\mathbf b=\mathbf 0)
\]
$\lambda>0$: same direction; \quad $\lambda<0$: opposite directions.
\end{theorem}

\begin{example}
(a) Show that $\cvv{2}{-4}{6}$ is parallel to $\cvv{-1}{2}{-3}$. \quad (b) Find $p$ and $q$ such that $\cvv{p}{6}{q}$ is parallel to $\cvv{1}{3}{-2}$.
\end{example}
\solution{\par\noindent\begin{minipage}[c]{0.62\linewidth}
\[
\begin{aligned}
\text{(a)}\quad\cvv{2}{-4}{6}&=-2\cvv{-1}{2}{-3}\\
\therefore\ \cvv{2}{-4}{6}&\parallel\cvv{-1}{2}{-3}\ \text{(shown)}\\[4pt]
\text{(b)}\quad\cvv{p}{6}{q}&=\lambda\cvv{1}{3}{-2}\\
6&=3\lambda\ \Rightarrow\ \lambda=2\\
\therefore\ p&=2,\quad q=-4
\end{aligned}
\]
\end{minipage}\hfill
\begin{minipage}[c]{0.34\linewidth}\centering
\begin{tikzpicture}[td,scale=0.3]
\axt{3}{-4}{4}{6}
\draw[ax](0,0,0)--(0,0,-3.6);
\node[dot,label={[L]below right:$O$}] at (0,0,0){};
\draw[vA](0,0,0)--(2,-4,6);\draw[vB](0,0,0)--(-1,2,-3);
\node[dot,label={[L]right:$(2,-4,6)$}] at (2,-4,6){};
\node[dot,label={[L]right:$(-1,2,-3)$}] at (-1,2,-3){};
\end{tikzpicture}
\end{minipage}}

\begin{theorem}{Unit vectors}
The unit vector in the direction of $\mathbf a$ is
\[
\hat{\mathbf a}=\frac{\mathbf a}{|\mathbf a|}
\]
\end{theorem}

\begin{example}
Let $\mathbf a=2\mathbf i-\mathbf j+2\mathbf k$. Find (a) the unit vector in the direction of $\mathbf a$, \quad (b) the vectors of magnitude $6$ parallel to $\mathbf a$.
\end{example}
\solution{\par\noindent\begin{minipage}[c]{0.62\linewidth}
\[
\begin{aligned}
\text{(a)}\quad|\mathbf a|&=\sqrt{2^2+(-1)^2+2^2}=\sqrt9=3\\
\hat{\mathbf a}&=\frac13\cvv{2}{-1}{2}\\[4pt]
\text{(b)}\quad\pm6\hat{\mathbf a}&=\pm\frac63\cvv{2}{-1}{2}=\pm\cvv{4}{-2}{4}
\end{aligned}
\]
\end{minipage}\hfill
\begin{minipage}[c]{0.34\linewidth}\centering
\begin{tikzpicture}[td,scale=0.9]
\draw[ax](0,0,0)--(3.1,0,0) node[right,L]{$x$};\draw[ax](0,0,0)--(0,0,3.2) node[above,L]{$z$};\draw[black!70](0,-1.5,0)--(0,2.4,0) node[above right,L]{$y$};
\node[dot,label={[L]below right:$O$}] at (0,0,0){};
\draw[vA](0,0,0)--(2,-1,2);
\draw[vR,line width=2.2pt](0,0,0)--(2/3,-1/3,2/3);
\node[L,blue!70!black,above left] at (1.5,-0.75,1.5){$\mathbf a$};
\node[green!45!black,below right] at (2/3,-1/3,2/3){$\hat{\mathbf a}$};
\node[dot,label={[L]right:$(2,-1,2)$}] at (2,-1,2){};
\end{tikzpicture}
\end{minipage}}

\clearpage
\subsection{Ratio Theorem}
\begin{theorem}{Ratio Theorem}
\noindent\begin{minipage}[c]{0.5\linewidth}
If $C$ lies on $AB$ with $AC:CB=\lambda:\mu$, then
\[
\vv{OC}=\frac{\mu\,\vv{OA}+\lambda\,\vv{OB}}{\lambda+\mu}
\]
i.e.\ $\mathbf c=\dfrac{\mu\mathbf a+\lambda\mathbf b}{\lambda+\mu}$.
\end{minipage}\hfill
\begin{minipage}[c]{0.46\linewidth}\centering
\begin{tikzpicture}[scale=1]
\coordinate (O) at (0,0);\coordinate (A) at (0.6,2.4);\coordinate (B) at (5,1.6);
\coordinate (C) at ($(A)!0.6!(B)$);
\draw[thick,black!70](A)--(B);
\draw[vA](O)--(A);\draw[vB](O)--(B);\draw[vR](O)--(C);
\node[dot,label=below left:$O$] at (O){};
\node[dot,label=above left:$A$] at (A){};
\node[dot,label=right:$B$] at (B){};
\node[dot,label=above:$C$] at (C){};
\node[blue!70!black,left] at ($(O)!0.5!(A)$){$\mathbf a$};
\node[red!70!black,below right] at ($(O)!0.5!(B)$){$\mathbf b$};
\node[green!45!black,above left=-1pt] at ($(O)!0.55!(C)$){$\mathbf c$};
\draw[decorate,decoration={brace,amplitude=4pt,raise=3pt}](A)--(C) node[midway,above=8pt,sloped]{$\lambda$};
\draw[decorate,decoration={brace,amplitude=4pt,raise=3pt}](C)--(B) node[midway,above=8pt,sloped]{$\mu$};
\end{tikzpicture}
\end{minipage}
\end{theorem}

\begin{example}
$\vv{OA}=\cvv{4}{1}{0}$ and $\vv{OB}=\cvv{1}{4}{6}$. $C$ lies on $AB$ with $AC:CB=2:1$. Find $\vv{OC}$.
\end{example}
\solution{\par\noindent\begin{minipage}[c]{0.36\linewidth}
\[
\begin{aligned}
\vv{OC}&=\frac{1\,\vv{OA}+2\,\vv{OB}}{2+1}\\
&=\frac13\left[\cvv{4}{1}{0}+\cvv{2}{8}{12}\right]\\
&=\frac13\cvv{6}{9}{12}\\
&=\cvv{2}{3}{4}
\end{aligned}
\]
\end{minipage}\hfill
\begin{minipage}[c]{0.3\linewidth}\centering
{\small Diagram you can draw:}\\[2pt]
\ratiosketch{$A(4,1,0)$}{$C\,?$}{$B(1,4,6)$}{2}{1}{0.667}
\end{minipage}\hfill
\begin{minipage}[c]{0.3\linewidth}\centering
\fcolorbox{black!30}{white}{\begin{minipage}{0.95\linewidth}\centering
{\scriptsize\color{black!50} Actual diagram: (you do not need to draw this)}\\[2pt]
\begin{tikzpicture}[td,scale=0.4,black!45]
\begin{scope}[every path/.style={black!40}]\draw[->,>=Stealth](0,0,0)--(5.5,0,0) node[right,L,black!50]{$x$};\draw[->,>=Stealth](0,-1,0)--(0,4.6,0) node[above right,L,black!50]{$y$};\draw[->,>=Stealth](0,0,0)--(0,0,6.8) node[above,L,black!50]{$z$};\end{scope}
\draw[black!45,thick](4,1,0)--(1,4,6);
\draw[->,>=Stealth,thick,black!45](0,0,0)--(4,1,0);\draw[->,>=Stealth,thick,black!45](0,0,0)--(1,4,6);\draw[->,>=Stealth,thick,black!60](0,0,0)--(2,3,4);
\node[dot,fill=black!45,label={[L,black!50]below right:$O$}] at (0,0,0){};
\node[dot,fill=black!45,label={[L,black!50]below:$A(4,1,0)$}] at (4,1,0){};
\node[dot,fill=black!45,label={[L,black!50]above right:$B(1,4,6)$}] at (1,4,6){};
\node[dot,fill=black!45,label={[L,black!50]right:$C(2,3,4)$}] at (2,3,4){};
\end{tikzpicture}
\end{minipage}}
\end{minipage}}

\begin{example}
$A$ has position vector $\cvv{2}{1}{-1}$ and $C$ has position vector $\cvv{5}{4}{5}$. $C$ lies on $AB$ with $AC:CB=3:2$. Find $\vv{OB}$.
\end{example}
\solution{\par\noindent\begin{minipage}[c]{0.36\linewidth}
\[
\begin{aligned}
\vv{OC}&=\frac{2\,\vv{OA}+3\,\vv{OB}}{3+2}\\
\therefore\ \vv{OB}&=\frac{5\,\vv{OC}-2\,\vv{OA}}{3}\\
&=\frac13\left[\cvv{25}{20}{25}-\cvv{4}{2}{-2}\right]\\
&=\frac13\cvv{21}{18}{27}\\
&=\cvv{7}{6}{9}
\end{aligned}
\]
\end{minipage}\hfill
\begin{minipage}[c]{0.3\linewidth}\centering
{\small Diagram you can draw:}\\[2pt]
\ratiosketch{$A(2,1,-1)$}{$C(5,4,5)$}{$B\,?$}{3}{2}{0.6}
\end{minipage}\hfill
\begin{minipage}[c]{0.3\linewidth}\centering
\fcolorbox{black!30}{white}{\begin{minipage}{0.95\linewidth}\centering
{\scriptsize\color{black!50} Actual diagram: (you do not need to draw this)}\\[2pt]
\begin{tikzpicture}[td,scale=0.27,black!45]
\begin{scope}[every path/.style={black!40}]\draw[->,>=Stealth](0,0,0)--(7.6,0,0) node[right,L,black!50]{$x$};\draw[->,>=Stealth](0,-1,0)--(0,6.6,0) node[above right,L,black!50]{$y$};\draw[->,>=Stealth](0,0,-1.6)--(0,0,9.8) node[above,L,black!50]{$z$};\end{scope}
\draw[black!45,thick](2,1,-1)--(7,6,9);
\draw[->,>=Stealth,thick,black!45](0,0,0)--(2,1,-1);\draw[->,>=Stealth,thick,black!60](0,0,0)--(5,4,5);\draw[->,>=Stealth,thick,black!45](0,0,0)--(7,6,9);
\node[dot,fill=black!45,label={[L,black!50]left:$O$}] at (0,0,0){};
\node[dot,fill=black!45,label={[L,black!50]below right:$A(2,1,-1)$}] at (2,1,-1){};
\node[dot,fill=black!45,label={[L,black!50]right:$C(5,4,5)$}] at (5,4,5){};
\node[dot,fill=black!45,label={[L,black!50]left:$B(7,6,9)$}] at (7,6,9){};
\end{tikzpicture}
\end{minipage}}
\end{minipage}}

\clearpage
\section{Scalar Product}
In this section, we learn about how we can multiply vectors together. There are two different ways we can do so, namely the Scalar Product and Vector Product. The Scalar and Vector products each have their own respective applications in the study of Vectors. We will learn about these applications more in depth later in this section.

\subsection{Scalar (Dot) Product}
\begin{definition}{Scalar (dot) product}
The scalar or dot product of two vectors $\mathbf a=a_1\mathbf i+a_2\mathbf j+a_3\mathbf k$ and $\mathbf b=b_1\mathbf i+b_2\mathbf j+b_3\mathbf k$ is defined as
\[
\mathbf a\cdot\mathbf b=\cvv{a_1}{a_2}{a_3}\cdot\cvv{b_1}{b_2}{b_3}=a_1b_1+a_2b_2+a_3b_3.
\]
\end{definition}

\begin{example}
Find $\cvv{1}{2}{3}\cdot\cvv{5}{-1}{3}$.
\end{example}
\solution{
\[
\begin{aligned}
\cvv{1}{2}{3}\cdot\cvv{5}{-1}{3}&=1(5)+2(-1)+3(3)\\
&=12
\end{aligned}
\]}


\needspace{15\baselineskip}
\subsection{Basic Properties of Scalar Product}
Taking reference to the problems in the Exercise Worksheet (Section 7, Scalar Product), we can see that the Scalar Product has properties that are similar to basic algebra multiplication. In general, we have the following properties of the Scalar Product.

\begin{theorem}{Laws of scalar product}
\begin{itemize}
\item $\mathbf a\cdot\mathbf b=\mathbf b\cdot\mathbf a$\\
{\small\color{black!50}$\mathbf a\cdot\mathbf b=22=\mathbf b\cdot\mathbf a$}
\item $\mathbf a\cdot(\mathbf b+\mathbf c)=\mathbf a\cdot\mathbf b+\mathbf a\cdot\mathbf c$,\quad $(\mathbf b+\mathbf c)\cdot\mathbf a=\mathbf b\cdot\mathbf a+\mathbf c\cdot\mathbf a$\\
{\small\color{black!50}$\mathbf a\cdot(\mathbf b+\mathbf c)=13=\mathbf a\cdot\mathbf b+\mathbf a\cdot\mathbf c$}
\item $\mathbf a\cdot(\lambda\mathbf b)=\lambda(\mathbf a\cdot\mathbf b)=(\lambda\mathbf a)\cdot\mathbf b,\quad\lambda\in\mathbb R$\\
{\small\color{black!50}$\mathbf a\cdot(2\mathbf b)=44=2(\mathbf a\cdot\mathbf b)=(2\mathbf a)\cdot\mathbf b$}
\item $|\mathbf a|^2=\mathbf a\cdot\mathbf a$\\
{\small\color{black!50}$\mathbf a\cdot\mathbf a=14=|\mathbf a|^2$}
\end{itemize}
{\small\color{black!50}where $\mathbf a=\cvv{1}{2}{3}$, $\mathbf b=\cvv{-3}{2}{7}$, $\mathbf c=\cvv{4}{-5}{-1}$.}
\end{theorem}



\needspace{15\baselineskip}
\subsection{Angle Between Two Vectors}
We now look at how we are able to use the Scalar Product to find the angle between two vectors. This is one of the applications of the Scalar Product. To be clear, when we say the angle $\theta$ between two vectors, we are referring to the non-reflex angle between them, where either both their start points coincide, or end points coincide.

\begin{center}
{\renewcommand{\arraystretch}{1.3}
\begin{tabular}{|*{4}{>{\centering\arraybackslash}m{0.2\textwidth}|}}
\hline
\multicolumn{2}{|c|}{\textbf{Correct angle}} & \multicolumn{2}{c|}{\textbf{Wrong angle}}\\
\hline
\vspace{4pt}{%
\begin{tikzpicture}[scale=0.85]
\draw[->,>=Stealth,thick](0,0)--(1.6,1.1) node[right]{$\mathbf a$};
\draw[->,>=Stealth,thick](0,0)--(1.6,-1.1) node[right]{$\mathbf b$};
\draw (0.5,0.34) arc[start angle=34.5,end angle=-34.5,radius=0.6];
\node at (0.9,0){$\theta$};
\end{tikzpicture}} &
\vspace{4pt}{%
\begin{tikzpicture}[scale=0.85]
\draw[->,>=Stealth,thick](1.6,1.1)--(0,0);
\draw[->,>=Stealth,thick](1.6,-1.1)--(0,0);
\node[above left] at (0.1,0.1){$\mathbf a$};\node[below left] at (0.1,-0.1){$\mathbf b$};
\draw (0.5,0.34) arc[start angle=34.5,end angle=-34.5,radius=0.6];
\node at (0.9,0){$\theta$};
\end{tikzpicture}} &
\vspace{4pt}{%
\begin{tikzpicture}[scale=0.85]
\draw[->,>=Stealth,thick](0,0)--(1.6,1.1) node[right]{$\mathbf a$};
\draw[->,>=Stealth,thick](0,0)--(1.6,-1.1) node[right]{$\mathbf b$};
\draw (0.5,0.34) arc[start angle=34.5,end angle=325.5,radius=0.6];
\node at (-0.25,0){$\theta$};
\end{tikzpicture}} &
\vspace{4pt}{%
\begin{tikzpicture}[scale=0.85]
\draw[->,>=Stealth,thick](1.6,1.1)--(0,0);
\draw[->,>=Stealth,thick](0,0)--(1.6,-1.1) node[below]{$\mathbf b$};
\node[above left] at (0.1,0.1){$\mathbf a$};
\draw (0.5,0.34) arc[start angle=34.5,end angle=-34.5,radius=0.6];
\node at (0.9,0){$\theta$};
\end{tikzpicture}}\vspace{4pt}\\
\hline
\end{tabular}}
\end{center}

In order to find this angle, we make use of the following formula.

\begin{definition}{(Definition of Scalar Product)}
The scalar product of two vectors $\mathbf a$ and $\mathbf b$ is given by
\[
\mathbf a\cdot\mathbf b=|\mathbf a||\mathbf b|\cos\theta,
\]
where $\theta$ is the angle between two vectors $\mathbf a$ and $\mathbf b$.
\end{definition}

\begin{info}
\textbf{Important Note:} For non-zero $\mathbf a$, $\mathbf b$,
\begin{itemize}
\item $\mathbf a\cdot\mathbf b>0\iff\theta$ is acute or zero ($0^\circ\le\theta<90^\circ$);
\item $\mathbf a\cdot\mathbf b<0\iff\theta$ is obtuse or $180^\circ$ ($90^\circ<\theta\le180^\circ$);
\item $\mathbf a\cdot\mathbf b=0\iff\theta=90^\circ$.
\end{itemize}
\end{info}

\begin{example}
Find the angle between the following vectors $\mathbf a=\cvv{1}{2}{3}$, $\mathbf b=\cvv{-2}{2}{1}$.
\end{example}
\solution{
\[
\begin{aligned}
\mathbf a\cdot\mathbf b&=|\mathbf a||\mathbf b|\cos\theta\\
\cvv{1}{2}{3}\cdot\cvv{-2}{2}{1}&=\left|\cvv{1}{2}{3}\right|\left|\cvv{-2}{2}{1}\right|\cos\theta\\
-2+4+3&=\sqrt{14}\sqrt9\cos\theta\\
\frac{5}{3\sqrt{14}}&=\cos\theta\\
\theta&=\cos^{-1}\left(\frac{5}{3\sqrt{14}}\right)\\
\theta&=63.5^\circ
\end{aligned}
\]}

Hence we can use Definition~13.4.2 to find the non-reflex angle $\theta$ between vectors $\mathbf a$ and $\mathbf b$:
\[
\cos\theta=\frac{\mathbf a\cdot\mathbf b}{|\mathbf a||\mathbf b|}.
\]


\needspace{12\baselineskip}
\subsection{Perpendicular Vectors}
We can see that from the previous results, that the scalar product of two non-zero vectors that are perpendicular to each other is zero.

\begin{theorem}{Perpendicular vectors}
\begin{center}
\begin{tikzpicture}[scale=0.9]
\draw[->,>=Stealth,thick](0,0)--(0,2.2) node[midway,left]{$\mathbf a$};
\draw[->,>=Stealth,thick](0,0)--(3,0) node[midway,below]{$\mathbf b$};
\draw(0.3,0)--(0.3,0.3)--(0,0.3);
\end{tikzpicture}
\end{center}
\begin{center}
If $\mathbf a$ and $\mathbf b$ are perpendicular vectors, then
\[
\mathbf a\cdot\mathbf b=0
\]
\end{center}
\end{theorem}

\textbf{Note:} This result can be used to test for two perpendicular vectors, provided that the two vectors are non-zero.

\begin{example}
Show that the vector $\mathbf a=\cvv{1}{2}{1}$ and $\mathbf b=\cvv{-1}{1}{-1}$ are perpendicular to each other.
\end{example}
\solution{
\[
\begin{aligned}
\mathbf a\cdot\mathbf b&=\cvv{1}{2}{1}\cdot\cvv{-1}{1}{-1}\\
&=-1+2-1\\
&=0
\end{aligned}
\]
Since $\mathbf a\cdot\mathbf b=0$, $\mathbf a$ \& $\mathbf b$ are perpendicular.}


\clearpage
\subsection{Length of Projection \& Projection Vector}
The projection vector of a vector $\mathbf a$ onto $\mathbf b$ is the vector component of $\mathbf a$ that is parallel to $\mathbf b$. The magnitude (or length) of this vector is known as the length of projection. To help you understand this better, imagine you are shining a torchlight from a point directly above $\mathbf a$. The length of projection is the length of the shadow that is cast on $\mathbf b$.

\begin{center}
\begin{tikzpicture}[scale=0.9]
\draw[->,>=Stealth,thick](0,0)--(4,3.4) node[midway,above left]{$\mathbf a$};
\draw[->,>=Stealth,thick](0,0)--(5.6,0);
\node[above] at (4.8,0){$\mathbf b$};
\draw[dashed](4,3.4)--(4,0);
\draw(3.75,0)--(3.75,0.25)--(4,0.25);
\node[left] at (0,0){$O$};\node[above right] at (4,3.4){$A$};\node[right] at (5.6,0){$B$};
\draw[|-|](0,-0.35)--(4,-0.35) node[midway,below]{length of projection};
\end{tikzpicture}
\end{center}

In order to find the length of projection and projection vector of $\mathbf a$ onto $\mathbf b$, we can make use of the following result.

\begin{theorem}{Projection of Vectors}
For any two non-zero vectors $\mathbf a$ and $\mathbf b$,
\[
\begin{aligned}
\text{length of projection of }\mathbf a\text{ onto }\mathbf b&=|\mathbf a\cdot\hat{\mathbf b}|,\\
\text{projection vector of }\mathbf a\text{ onto }\mathbf b&=(\mathbf a\cdot\hat{\mathbf b})\hat{\mathbf b}.
\end{aligned}
\]
\end{theorem}

\begin{example}
The vectors $\mathbf a$ and $\mathbf b$ are given by
\[
\mathbf a=\cvv{1}{2}{3},\quad\text{and}\quad\mathbf b=\cvv{-1}{-2}{2}.
\]
Find the length of projection and the projection vector of the vector $\mathbf a$ onto $\mathbf b$.
\end{example}
\solution{\par\noindent\begin{minipage}[c]{0.6\linewidth}
\[
\begin{aligned}
\text{Length of projection}&=|\mathbf a\cdot\hat{\mathbf b}|\\
\hat{\mathbf b}&=\frac13\cvv{-1}{-2}{2}\\
\therefore\ \text{length}&=\left|\cvv{1}{2}{3}\cdot\frac13\cvv{-1}{-2}{2}\right|\\
&=\frac13\left|\cvv{1}{2}{3}\cdot\cvv{-1}{-2}{2}\right|\\
&=\frac13\text{ units}
\end{aligned}
\]
\end{minipage}\hfill
\begin{minipage}[c]{0.36\linewidth}\centering
\begin{tikzpicture}[scale=0.8]
\draw[->,>=Stealth,thick](0,0)--(3,2.6) node[midway,above left]{$\mathbf a$};
\draw[->,>=Stealth,thick](0,0)--(3.8,0) node[below]{$\mathbf b$};
\draw[dashed](3,2.6)--(3,0);
\draw(2.75,0)--(2.75,0.25)--(3,0.25);
\node[left] at (0,0){$O$};
\end{tikzpicture}
\end{minipage}
\par
$\therefore$ The length of projection of the vector $\mathbf a$ onto $\mathbf b$ is $\dfrac13$ units. The projection vector is
\[
\frac13\left(\frac13\cvv{-1}{-2}{2}\right)=\frac19\cvv{-1}{-2}{2}
\]}


\clearpage
\section{Vector Product}
In this section, we learn how to multiply vectors together using the Vector Product. The Vector Product is slightly more difficult to apply, in comparison to the Scalar Product.

\begin{definition}{Vector (cross) product}
The vector or cross product of two vectors $\mathbf a=a_1\mathbf i+a_2\mathbf j+a_3\mathbf k$ and $\mathbf b=b_1\mathbf i+b_2\mathbf j+b_3\mathbf k$ (denoted by $\mathbf a\times\mathbf b$) is
\[
\mathbf a\times\mathbf b=\cvv{a_1}{a_2}{a_3}\times\cvv{b_1}{b_2}{b_3}=\cvv{a_2b_3-a_3b_2}{a_3b_1-a_1b_3}{a_1b_2-a_2b_1}
\]
\end{definition}

\begin{example}
Given that $\mathbf a=\cvv{1}{2}{3}$, $\mathbf b=\cvv{-1}{3}{-2}$, find $\mathbf a\times\mathbf b$.
\end{example}
\solution{
\[
\begin{aligned}
\cvv{1}{2}{3}\times\cvv{-1}{3}{-2}&=\cvv{2(-2)-3(3)}{3(-1)-1(-2)}{1(3)-2(-1)}\\
\therefore\ \mathbf a\times\mathbf b&=\cvv{-13}{-1}{5}
\end{aligned}
\]}


\needspace{12\baselineskip}
\subsection{Basic Properties of Vector Product}
Taking reference to the problems in the Exercise Worksheet (Section 8, Vector Product), we can see that the Vector Product has properties that are similar to that of scalar product, and basic algebra. While some laws do not hold true for the vector product, others still hold true.

\begin{theorem}{Laws of vector product}
\begin{itemize}
\item $\mathbf a\times\mathbf b=-(\mathbf b\times\mathbf a)$\\
{\small\color{black!50}$\cvv{1}{2}{3}\times\cvv{2}{-1}{4}=\cvv{11}{2}{-5}=-\left(\cvv{2}{-1}{4}\times\cvv{1}{2}{3}\right)$}
\item $\mathbf a\times(\mathbf b+\mathbf c)=\mathbf a\times\mathbf b+\mathbf a\times\mathbf c$\\
{\small\color{black!50}$\cvv{1}{2}{3}\times\cvv{1}{-3}{6}=\cvv{21}{-3}{-5}=\cvv{1}{2}{3}\times\cvv{-3}{2}{7}+\cvv{1}{2}{3}\times\cvv{4}{-5}{-1}$}
\item $\lambda(\mathbf a\times\mathbf b)=(\lambda\mathbf a)\times\mathbf b=\mathbf a\times(\lambda\mathbf b),\quad\lambda\in\mathbb R$\\
{\small\color{black!50}$2\left(\cvv{1}{2}{3}\times\cvv{-3}{2}{7}\right)=\cvv{16}{-32}{16}=\cvv{2}{4}{6}\times\cvv{-3}{2}{7}=\cvv{1}{2}{3}\times\cvv{-6}{4}{14}$}
\item $\mathbf a\times\mathbf a=\mathbf 0$\\
{\small\color{black!50}$\cvv{1}{2}{3}\times\cvv{1}{2}{3}=\mathbf 0$}
\end{itemize}
\end{theorem}

\needspace{12\baselineskip}
\subsection{Resultant Vector of the Vector Product}
When we apply the Vector Product to two vectors $\mathbf a$ and $\mathbf b$, we get a resultant vector $\mathbf a\times\mathbf b$. Is there a significant relationship of the vector $\mathbf a\times\mathbf b$ to the vectors $\mathbf a$ and $\mathbf b$? There is, in fact a relationship between these three vectors.

\begin{center}
\begin{tikzpicture}[td,scale=1.3]
\draw[->,>=Stealth,thick](0,0,0)--(2.2,0,0) node[midway,above]{$\mathbf b$};
\draw[->,>=Stealth,thick](0,0,0)--(0,-1.8,0) node[left]{$\mathbf a$};
\draw[->,>=Stealth,thick](0,0,0)--(0,0,1.8) node[right]{$\mathbf a\times\mathbf b$};
\draw(0.22,0,0)--(0.22,0,0.22)--(0,0,0.22);
\draw(0,-0.25,0)--(0,-0.25,0.22)--(0,0,0.22);
\end{tikzpicture}
\end{center}

\textbf{Note:} The vector $\mathbf b\times\mathbf a$ would also be perpendicular to the vectors $\mathbf a$ and $\mathbf b$, but in the opposite direction as $\mathbf a\times\mathbf b$, because $\mathbf b\times\mathbf a=-\mathbf a\times\mathbf b$.

The vectors $\mathbf a\times\mathbf b$ and $\mathbf b\times\mathbf a$ are known as normal vectors to the plane in which the vectors $\mathbf a$ and $\mathbf b$ lies in.

\begin{example}
Given that $\mathbf a=3\mathbf i-5\mathbf j+6\mathbf k$ and $\mathbf b=3\mathbf i+3\mathbf j-2\mathbf k$, find a vector perpendicular to both $\mathbf a$ and $\mathbf b$.
\end{example}
\solution{A vector perpendicular to $\mathbf a$ and $\mathbf b$ is
\[
\begin{aligned}
\mathbf a\times\mathbf b&=\cvv{3}{-5}{6}\times\cvv{3}{3}{-2}\\
&=\cvv{(-5)(-2)-(3)(6)}{(6)(3)-(-2)(3)}{(3)(3)-(3)(-5)}\\
&=\cvv{-8}{24}{24}
\end{aligned}
\]}

\begin{question}
\textbf{Thinking question:} \textit{Can you think of other vectors that are perpendicular to both $\mathbf a$ and $\mathbf b$?}
\end{question}

\clearpage
\subsection{Parallel Vectors}
We now consider the cross product of two vectors parallel to each other. Suppose $\mathbf a$ and $\mathbf b$ are two non-zero vectors that are parallel to each other.

\begin{theorem}{Parallel vectors}
\begin{center}
\begin{tikzpicture}[scale=0.9]
\draw[->,>=Stealth,thick](0,0)--(1.4,1) node[midway,above left]{$\mathbf a$};
\draw[->,>=Stealth,thick](1.2,0)--(4,2) node[midway,below right]{$\mathbf b$};
\end{tikzpicture}
\[
\begin{aligned}
&\text{If }\mathbf a/\!/\mathbf b,\\
&\implies\mathbf a\times\mathbf b=\mathbf 0
\end{aligned}
\]
The cross product of any two non-zero, parallel vectors will be equal to $\mathbf 0$.
\end{center}
\end{theorem}

We consider an explicit example of the above property.

\begin{example}
Find $\cvv{1}{-2}{1}\times\cvv{-3}{6}{-3}$.
\end{example}
\solution{
\[
\begin{aligned}
\cvv{1}{-2}{1}\times\cvv{-3}{6}{-3}&=\cvv{6-6}{-3-(-3)}{6-6}\\
&=\cvv{0}{0}{0}
\end{aligned}
\]
Since $\cvv{1}{-2}{1}\times\cvv{-3}{6}{-3}=\cvv{0}{0}{0}$, the vectors $\cvv{1}{-2}{1}$ and $\cvv{-3}{6}{-3}$ are parallel to each other.}

\begin{example}
Show that $\mathbf a=\cvv{1}{-2}{1}$ and $\mathbf b=\cvv{2}{-4}{2}$ are parallel.
\end{example}
\solution{\par\noindent\begin{minipage}[t]{0.55\linewidth}
Method 1:
\[
\begin{aligned}
\cvv{1}{-2}{1}\times\cvv{2}{-4}{2}&=\cvv{(-2)(2)-(1)(-4)}{(1)(2)-(1)(2)}{(1)(-4)-(-2)(2)}\\
&=\cvv{0}{0}{0}
\end{aligned}
\]
Since $\mathbf a$ and $\mathbf b$ are non-zero vectors, $\mathbf a$ is parallel to $\mathbf b$.
\end{minipage}\hfill
\begin{minipage}[t]{0.4\linewidth}
Method 2:
\[
\begin{aligned}
\mathbf b&=\cvv{2}{-4}{2}\\
&=2\cvv{1}{-2}{1}\\
&=2\mathbf a
\end{aligned}
\]
Since $\mathbf b$ is a scalar multiple of $\mathbf a$, $\mathbf a$ and $\mathbf b$ are parallel vectors.
\end{minipage}}



\subsection{Vector Product Formula}
Similar to the dot product formula $[\mathbf a\cdot\mathbf b=|\mathbf a||\mathbf b|\cos\theta]$ the cross product has a formula as well.

\begin{theorem}{Vector product formula}
\[
|\mathbf a\times\mathbf b|=|\mathbf a||\mathbf b|\sin\theta
\]
Where $\theta$ is the angle between the vectors $\mathbf a$ and $\mathbf b$.
\end{theorem}

Note that we have to include the modulus sign as $|\mathbf a\times\mathbf b|$ instead of $\mathbf a\times\mathbf b$, because $\mathbf a\times\mathbf b$ would return us a vector, not a scalar.

\begin{example}
Find the acute angle between the two vectors $\mathbf a=\cvv{1}{2}{1}$ and $\mathbf b=\cvv{-1}{3}{2}$ by using vector product.
\end{example}
\solution{
\[
\begin{aligned}
\left|\cvv{1}{2}{1}\times\cvv{-1}{3}{2}\right|&=\left|\cvv{1}{2}{1}\right|\left|\cvv{-1}{3}{2}\right|\sin\theta\\
\left|\cvv{4-3}{-1-2}{3-(-2)}\right|&=\sqrt6\sqrt{14}\sin\theta\\
\left|\cvv{1}{-3}{5}\right|&=\sqrt6\sqrt{14}\sin\theta\\
\frac{\sqrt{35}}{\sqrt6\sqrt{14}}&=\sin\theta\\
\theta&=40.2^\circ
\end{aligned}
\]}




\clearpage
\subsection{Application of Vector Product: Perpendicular Distance}
In this section, we deal with the perpendicular component of vectors. This concept is similar to length of projection of vectors in the previous section.

\begin{center}
\begin{tikzpicture}[scale=0.9]
\draw[->,>=Stealth,thick](0,0)--(4,3) node[midway,above left]{$\mathbf a$};
\draw[->,>=Stealth,thick](0,0)--(6,0);
\node[below] at (2,0){$\mathbf b$};
\draw[dashed](4,3)--(4,0);
\draw(3.75,0)--(3.75,0.25)--(4,0.25);
\node[left] at (0,0){$O$};\node[above] at (4,3){$A$};\node[right] at (6,0){$B$};\node[below] at (4,0){$F$};
\end{tikzpicture}

The length $OF$ is the length of projection of $\mathbf a$ on $\mathbf b$.\\
The length $AF$ is the perpendicular distance from $A$ to the line joining $O$ and $B$.
\end{center}

\begin{theorem}{Perpendicular component}
For any two non-zero vectors $\mathbf a$ and $\mathbf b$,
\begin{enumerate}[label=\alph*)]
\item Length of vector component of $\mathbf a$ perpendicular to $\mathbf b$ (Length $AF$) $=|\mathbf a\times\hat{\mathbf b}|$.
\item Vector component of $\mathbf a$ perpendicular to $\mathbf b=\mathbf a-(\mathbf a\cdot\hat{\mathbf b})\hat{\mathbf b}$.
\end{enumerate}
\end{theorem}

\textbf{Note:} You can think of this concept as the ``cross product version'' of the length of projection.

\begin{example}
Relative to an origin $O$, two points $A$ and $B$ have position vectors given by $\mathbf a=\mathbf i+2\mathbf j+2\mathbf k$ and $\mathbf b=2\mathbf i-\mathbf j+3\mathbf k$. Find the shortest distance from $B$ to the line joining $O$ and $A$.
\end{example}
\solution{\par\noindent\begin{minipage}[c]{0.55\linewidth}
\[
\begin{aligned}
\text{Required distance}&=|\mathbf b\times\hat{\mathbf a}|\\
&=\left|\cvv{2}{-1}{3}\times\frac13\cvv{1}{2}{2}\right|\\
&=\frac13\left|\cvv{-8}{-1}{5}\right|\\
&=\frac13\sqrt{90}\\
&=\sqrt{10}\text{ units}
\end{aligned}
\]
\end{minipage}\hfill
\begin{minipage}[c]{0.42\linewidth}\centering
\begin{tikzpicture}[scale=0.8]
\draw[thick](0,0)--(5,0);
\draw[thick](0,0)--(3.2,2.8);
\draw[dashed](3.2,2.8)--(3.2,0);
\draw(2.95,0)--(2.95,0.25)--(3.2,0.25);
\node[left] at (0,0){$O$};\node[right] at (5,0){$A$};\node[above] at (3.2,2.8){$B$};\node[below] at (3.2,0){$F$};
\node[left] at (1.4,1.7){$\cvv{2}{-1}{3}$};\node[below] at (4.5,-0.1){$\cvv{1}{2}{2}$};
\end{tikzpicture}
\end{minipage}}


\clearpage

\clearpage
\subsection{Areas of Triangles \& Parallelograms}
Another application of the Vector Product is to find areas of triangles and parallelograms.

\begin{theorem}{Areas}
\begin{center}
\begin{tikzpicture}[scale=0.85]
\draw[->,>=Stealth,thick](0,0)--(1.5,2) node[midway,left]{$\mathbf a$};
\draw[->,>=Stealth,thick](0,0)--(4,0) node[midway,below]{$\mathbf b$};
\draw[thick](1.5,2)--(4,0);
\draw[dashed](1.5,2)--(5.5,2)--(4,0);
\node[left] at (0,0){$O$};\node[above] at (1.5,2){$A$};\node[below right] at (4,0){$B$};
\end{tikzpicture}
\[
\begin{aligned}
\text{Area of triangle }OAB&=\frac12|\mathbf a\times\mathbf b|\\
\text{Area of parallelogram with sides }\mathbf a\text{ and }\mathbf b&=|\mathbf a\times\mathbf b|
\end{aligned}
\]
\end{center}
\end{theorem}

\textbf{Note:} We can use the area of triangle formula, $\frac12ab\sin C$ to see why this formula holds true.

\begin{example}
Relative to an origin $O$, two points $A$ and $B$ have position vectors given by $\mathbf a=\mathbf i+2\mathbf j+2\mathbf k$ and $\mathbf b=2\mathbf i-\mathbf j+3\mathbf k$. Find the area of the triangle $OAB$.
\end{example}
\solution{\par\noindent\begin{minipage}[c]{0.55\linewidth}
\[
\begin{aligned}
\text{Required area}&=\frac12|\vv{OA}\times\vv{OB}|\\
&=\frac12\left|\cvv{1}{2}{2}\times\cvv{2}{-1}{3}\right|\\
&=\frac12\left|\cvv{8}{1}{-5}\right|\\
&=\frac12\sqrt{90}\\
&=\frac32\sqrt{10}\text{ units}^2
\end{aligned}
\]
\end{minipage}\hfill
\begin{minipage}[c]{0.42\linewidth}\centering
\begin{tikzpicture}[scale=0.75]
\draw[thick](0,0)--(3.5,2.6)--(4.5,0)--cycle;
\node[left] at (0,0){$O$};
\node[right] at (3.5,2.6){$A\cvv{1}{2}{2}$};
\node[right] at (4.5,0){$B\cvv{2}{-1}{3}$};
\end{tikzpicture}
\end{minipage}}

We can now extend this concept to find areas of triangles and parallelograms without the origin $O$ as a vertex.

\begin{center}
\begin{tikzpicture}[scale=0.75]
\draw[thick](0,0)--(4.5,0)--(1.8,1.8)--cycle;
\node[left] at (0,0){$A$};\node[right] at (4.5,0){$B$};\node[above] at (1.8,1.8){$C$};
\end{tikzpicture}
\hspace{1.5cm}
\begin{tikzpicture}[scale=0.75]
\draw[thick](0,0)--(4.5,0)--(6,2.5)--(1.5,2.5)--cycle;
\node[left] at (0,0){$A$};\node[right] at (4.5,0){$B$};\node[above] at (6,2.5){$C$};\node[above] at (1.5,2.5){$D$};
\end{tikzpicture}
\end{center}

Consider the above triangle and parallelogram. The area for the respective figures can be obtained from the following result.
\begin{theorem}{Areas without $O$ as a vertex}
\[
\begin{aligned}
\text{Area of triangle }ABC&=\frac12|\vv{AB}\times\vv{AC}|\\
\text{Area of parallelogram }ABCD&=|\vv{AB}\times\vv{AD}|
\end{aligned}
\]
\end{theorem}

\clearpage
\begin{example}
The position vectors $A$, $B$ and $C$ are $\mathbf a=\mathbf i+2\mathbf j+\mathbf k$, $\mathbf b=\mathbf i+3\mathbf k$ and $\mathbf c=-\mathbf i+2\mathbf j+7\mathbf k$. Find the area of the triangle $ABC$.
\end{example}
\solution{\par\noindent\begin{minipage}[c]{0.58\linewidth}
\[
\begin{aligned}
\text{Required Area}&=\frac12|\vv{AB}\times\vv{AC}|\\
\vv{AB}&=\vv{OB}-\vv{OA}=\cvv{1}{0}{3}-\cvv{1}{2}{1}=\cvv{0}{-2}{2}\\
\vv{AC}&=\vv{OC}-\vv{OA}=\cvv{-1}{2}{7}-\cvv{1}{2}{1}=\cvv{-2}{0}{6}\\
\therefore\ \text{Area}&=\frac12\left|\cvv{0}{-2}{2}\times\cvv{-2}{0}{6}\right|\\
&=\frac12\left|\cvv{-12}{-4}{-4}\right|\\
&=\frac{\sqrt{176}}{2}\\
&=2\sqrt{11}\text{ units}^2
\end{aligned}
\]
\end{minipage}\hfill
\begin{minipage}[c]{0.38\linewidth}\centering
\begin{tikzpicture}[scale=0.8]
\draw[thick](0,0)--(2,1.6)--(4,0)--cycle;
\node[left] at (0,0){$A$};\node[above] at (2,1.6){$B$};\node[right] at (4,0){$C$};
\end{tikzpicture}
\end{minipage}}



\clearpage
\subsection{Summary of Scalar \& Vector Products}
%% ---- Direct copy of Pure Math Summary §15.2 table (ch15_notes.tex) ----
%% Source: /workspace/flat/summaries/systems-of-linear-equations/hw/ch15_notes.tex
%% Reused: tabular colspec m{2.4cm}|m{6.2cm}|m{6.2cm}, arraystretch 1.25, \small,
%%   \Xhline{2pt}, \makecell rows, \sidenote grey notes, \topdisplay header,
%%   tikz ar/pt styles, geometrical-meaning nested tabular + three diagrams,
%%   Note row, post-table info + Useful identity.
%% Adaptations (unavoidable name clashes with Vectors I notes macros):
%%   \vv → \mathbf (notes use \vv = overrightarrow); \cv → 3-arg pmatrix
%%   (notes use 2-arg); local \sidenote/\topdisplay/\hl/\vab; PMS-style
%%   theorem (no auto-number) for Useful identity.
\begingroup
\renewcommand{\vv}[1]{\mathbf{#1}}%
\renewcommand{\cv}[3]{\begin{pmatrix}#1\\#2\\#3\end{pmatrix}}%
\newcommand{\pmsidenote}[1]{{\small\textcolor{gray}{#1}}}%
\let\sidenote\pmsidenote
\newcommand{\topdisplay}[1]{\par\noindent\hfill$\displaystyle #1$\hfill\null\par\smallskip}%
\renewcommand{\hl}[1]{#1}%
\providecommand{\vab}[1]{\overrightarrow{#1}}%
\tikzset{pt/.style={circle,fill,inner sep=1.2pt},ar/.style={-{Stealth},thick}}%
% PMS-style theorem (title only; Vectors notes theorem auto-numbers)
\renewenvironment{theorem}[2][]{%
  \begin{tcolorbox}[enhanced, before skip=2mm, after skip=2mm,
    colback=black!5!white,colframe=black,fonttitle=\bfseries,title={#2},#1]}{\end{tcolorbox}}%
\topdisplay{\vv a=a_1\vv i+a_2\vv j+a_3\vv k,\qquad \vv b=b_1\vv i+b_2\vv j+b_3\vv k,\qquad\theta=\text{non-reflex }\angle AOB}
\begin{center}\renewcommand{\arraystretch}{1.25}\small
\begin{tabular}{|>{\centering\arraybackslash}m{2.4cm}|>{\centering\arraybackslash}m{6.2cm}|>{\centering\arraybackslash}m{6.2cm}|}\hline
 & \textbf{Scalar Product (Dot)} & \textbf{Vector Product (Cross)}\\\Xhline{2pt}
Definition & $\vv a\cdot\vv b=a_1b_1+a_2b_2+a_3b_3$ & $\vv a\times\vv b=\cv{a_2b_3-a_3b_2}{-(a_1b_3-a_3b_1)}{a_1b_2-a_2b_1}$\\\hline
Result & scalar & vector, $\perp$ to both $\vv a$ and $\vv b$\\\hline
Formula & $\vv a\cdot\vv b=|\vv a||\vv b|\cos\theta$ & $|\vv a\times\vv b|=|\vv a||\vv b|\sin\theta$\\\hline
Angle & $\cos\theta=\dfrac{\vv a\cdot\vv b}{|\vv a||\vv b|}$\par\sidenote{\hl{$\vv a\cdot\vv b>0$ acute or zero, $<0$ obtuse or $180^\circ$, $=0$ right angle}} & $\sin\theta=\dfrac{|\vv a\times\vv b|}{|\vv a||\vv b|}$\par
\sidenote{this method only finds the basic angle}\par
\sidenote{possible, but usually undesirable}\\\hline
Properties &
\makecell{$\vv a\cdot\vv b=\vv b\cdot\vv a$\\ $\vv a\cdot(\vv b+\vv c)=\vv a\cdot\vv b+\vv a\cdot\vv c$\\ $\lambda(\vv a\cdot\vv b)=(\lambda\vv a)\cdot\vv b=\vv a\cdot(\lambda\vv b)$\\ $\vv a\cdot\vv a=|\vv a|^{2}$} &
\makecell{$\vv a\times\vv b=-\vv b\times\vv a$\\ $\vv a\times(\vv b+\vv c)=\vv a\times\vv b+\vv a\times\vv c$\\ $\lambda(\vv a\times\vv b)=(\lambda\vv a)\times\vv b=\vv a\times(\lambda\vv b)$\\ $\vv a\times\vv a=\vv 0$}\\\hline
Test & $\vv a\cdot\vv b=0\iff\vv a\perp\vv b$ & $\vv a\times\vv b=\vv 0\iff\vv a\,/\!/\,\vv b$\\
 & \multicolumn{2}{c|}{\sidenote{\hl{(non-zero $\vv a$, $\vv b$)}}}\\\hline
Applications &
\makecell{Length of projection $=|\vv a\cdot\hat{\vv b}|$\\ Projection vector $=(\vv a\cdot\hat{\vv b})\hat{\vv b}$} &
\makecell{Perpendicular distance $=|\vv a\times\hat{\vv b}|$\\ $\text{Area of }\triangle OAB=\tfrac12|\vv a\times\vv b|$\\ Area of parallelogram $=|\vv a\times\vv b|$}\\\hline
Geometrical meaning &
\begin{tikzpicture}[scale=0.5]
\draw[ar](0,0)--(3.4,0)node[right]{$\vv b$};\draw[ar](0,0)--(2.2,1.6)node[above]{$\vv a$};
\draw[dashed](2.2,1.6)--(2.2,0);\draw[very thick,red](0,0)--(2.2,0);\node[below,red,font=\footnotesize]at(1.1,0){$|\vv a\cdot\hat{\vv b}|$};
\end{tikzpicture}\par
$|\vv a\cdot\hat{\vv b}|$ = length of projection of $\vv a$ onto $\vv b$
&
\scriptsize\renewcommand{\arraystretch}{1}\setlength{\tabcolsep}{1.5pt}\begin{tabular}{@{}>{\centering\arraybackslash}m{1.9cm}>{\centering\arraybackslash}m{1.9cm}>{\centering\arraybackslash}m{1.9cm}@{}}
\begin{tikzpicture}[scale=0.36]
\draw[dashed](-0.6,0)--(4.4,0);
\draw[ar](0,0)--(3.2,0)node[below]{$\vv b$};\draw[ar](0,0)--(1.4,2)node[above left]{$\vv a$};
\draw[very thick,red](1.4,2)--(1.4,0);\draw(1.4,0)++(0.25,0)--++(0,0.25)--++(-0.25,0);
\node[below left]at(0,0){$O$};
\end{tikzpicture} &
\begin{tikzpicture}[scale=0.36]
\draw[fill=blue!8](0,0)--(3,0)--(4.4,2)--(1.4,2)--cycle;
\draw[ar](0,0)--(3,0)node[below]{$\vv b$};\draw[ar](0,0)--(1.4,2)node[above left]{$\vv a$};
\node[below left]at(0,0){$O$};
\end{tikzpicture} &
\begin{tikzpicture}[scale=0.36]
\draw[dashed](3,0)--(4.4,2)--(1.4,2);
\draw[fill=blue!15](0,0)--(3,0)--(1.4,2)--cycle;
\draw[ar](0,0)--(3,0)node[below]{$B$};\draw[ar](0,0)--(1.4,2)node[above left]{$A$};
\node[below left]at(0,0){$O$};
\end{tikzpicture}\\[1mm]
perpendicular distance from $A$ to line along $\vv b$\par{\color{red}$|\vv a\times\hat{\vv b}|$} &
area of parallelogram, sides $\vv a$, $\vv b$\par$|\vv a\times\vv b|$ &
$2\times$ area of $\triangle OAB$\par$|\vv a\times\vv b|$
\end{tabular}\\\hline
Note & \sidenote{$\vv a\cdot\vv b\cdot\vv c$ meaningless; no cancellation law} & \sidenote{no cancellation law}\par\sidenote{$(\vv a\times\vv b)\times\vv c$ is a vector, but triple cross products are usually not tested}\\\hline
\end{tabular}\end{center}
\begin{info}[title=Note, before skip=2mm, after skip=2mm]
Projection of $\vv a$ onto $\vv b$ is not generally equal to projection of $\vv b$ onto $\vv a$. Perpendicular distance from $A$ to $OB$ is not generally equal to that from $B$ to $OA$. Triangles without $O$ as a vertex: $\tfrac12|\vab{AB}\times\vab{AC}|$.
\end{info}
\begin{theorem}[before skip=0pt, after skip=0pt]{Useful identity}
For any vectors $\vv a$, $\vv b$:\quad $(\vv a\times\vv b)\cdot\vv a=0$\quad and\quad $(\vv a\times\vv b)\cdot\vv b=0$\par
$\because\ \vv a\times\vv b$ is perpendicular to $\vv a$ and $\vv b$.
\end{theorem}
\endgroup

\end{document}
